\documentclass[11pt]{article}
\usepackage[T1]{fontenc}
\usepackage[utf8]{inputenc}
\usepackage{lmodern,amsmath,amssymb,amsthm,mathtools,mathrsfs,geometry,microtype,graphicx,booktabs,longtable,array}
\usepackage[hidelinks]{hyperref}
\geometry{margin=27mm}
\setlength{\emergencystretch}{2em}
\raggedbottom
\numberwithin{equation}{section}
\newtheorem{theorem}{Theorem}[section]
\newtheorem{lemma}[theorem]{Lemma}
\newtheorem{proposition}[theorem]{Proposition}
\newtheorem{corollary}[theorem]{Corollary}
\theoremstyle{definition}
\newtheorem{definition}[theorem]{Definition}
\newtheorem{remark}[theorem]{Remark}
\newcommand{\C}{\mathbb C}
\newcommand{\R}{\mathbb R}
\newcommand{\Z}{\mathbb Z}
\newcommand{\T}{\mathbb T}
\newcommand{\dd}{\,\mathrm d}
\newcommand{\ii}{\mathrm i}
\newcommand{\eps}{\varepsilon}
\newcommand{\Pf}{\operatorname{Pf}}
\newcommand{\Res}{\operatorname{Res}}
\newcommand{\dist}{\operatorname{dist}}
\newcommand{\cS}{\mathcal S}
\newcommand{\cA}{\mathscr A}
\newcommand{\OH}{\Omega_H}
\newcommand{\OL}{\Omega_L}
\newcommand{\XH}{\mathcal X_H}
\newcommand{\XL}{\mathcal X_L}
\newcommand{\norm}[1]{\left\lVert#1\right\rVert}
\title{Natural boundaries of the full bulk susceptibility\newline
for the fixed-ratio anisotropic square-lattice Ising model}
\author{Mathematical research campaign ANISO-ISING-BULK-NB-ALL-LAMBDA-LT-HT}
\date{19 October 2026\newline Research version 1.0: internally checked analytic proof}
\begin{document}
\maketitle
\begin{abstract}
For every fixed finite positive coupling ratio, on both physical
temperature sides, we establish explicit natural boundaries for the
complete bulk magnetic susceptibility. The independently defined
spectral set is first determined exactly on the one-complex-variable
temperature slice. Its high-temperature component is a bounded lens and
its low-temperature component is a simply connected epigraph. Regular
positive spectral edges, and both signs for the even-particle low-temperature
sum, admit nonzero first singular amplitudes with a Gaussian bound on the
entire differentiated particle tail. Negative high-temperature edges have
no exact odd endpoint resonance. There the argument uses an exact
permutation-equivariant pole normal form, a dimension-uniform beta-period
estimate, and dominance of a single particle index in selected high Taylor
coefficients. The isotropic case is separately matched, without a limiting
argument, to the reused normalized isotropic theorems at their exact
coefficient scope. The physical normalization, every lattice site, and the
pure-phase low-temperature magnetization subtraction are retained. The
proof uses the cited real-temperature spin form factors and explicitly
identified isotropic analytic results; no anisotropic natural-boundary
claim is an input. Three separate internal adversarial reconstructions
and their reproducible diagnostics accompany the manuscript. They do not
constitute external peer review or a new formal-kernel validation.
\end{abstract}
\tableofcontents

% BEGIN embedded section: sections/01_scope_physical.tex
\section{Fixed-temperature slice and the physical observable}\label{sec:physical}

Fix $J_x>0$ and $\lambda=J_y/J_x\in(0,\infty)$. Throughout,
\begin{equation}\label{eq:variables}
 z=2\beta J_x,\qquad S=\sinh z,\quad T=\sinh(\lambda z),\quad
 A=\cosh z\cosh(\lambda z),\quad p=ST.
\end{equation}
The ratio $\lambda$ is fixed before any complex continuation; in particular,
the generally varying ratio $T/S$ is not the fixed coupling ratio.
There is a unique $z_c>0$ with $\sinh z_c\sinh(\lambda z_c)=1$.
Define strictly increasing functions on $[0,\infty)$ by
\begin{equation}\label{eq:physicalPQ}
 P_\lambda(x)=\cosh x\tanh(\lambda x),\qquad
 Q_\lambda(x)=\cosh(\lambda x)\tanh x.
\end{equation}
Their values at $z_c$ are both one. Set
\begin{align}
 Y_\lambda&=\frac{\pi}{2\max(1,\lambda)},\nonumber\\
 H_\lambda(y)&=\min\{P_\lambda^{-1}(\cos y),
                      Q_\lambda^{-1}(\cos(\lambda y))\},
                      &&|y|<Y_\lambda,\label{eq:Hdef}\\
 L_\lambda(y)&=\max\{P_\lambda^{-1}(|\cos y|),
                      Q_\lambda^{-1}(|\cos(\lambda y)|)\},
                      &&y\in\R.\label{eq:Ldef}
\end{align}
The domains of interest are specified independently of susceptibility:
\begin{equation}\label{eq:domains}
 \OH=\{x+\ii y:|y|<Y_\lambda,\ |x|<H_\lambda(y)\},\qquad
 \OL=\{x+\ii y:x>L_\lambda(y)\}.
\end{equation}
The value of $H_\lambda$ at each endpoint is its continuous limit zero.

\begin{theorem}[Full bulk susceptibility at every positive ratio]\label{thm:aniso}
Let $0<\lambda<\infty$. The high-temperature physical
bulk susceptibility germ extends holomorphically to $\OH$, and
$\partial\OH$ is its natural boundary. The low-temperature physical
pure-phase connected bulk susceptibility germ extends holomorphically to
$\OL$, and the entire graph $\partial\OL=\{L_\lambda(y)+\ii y:y\in\R\}$
is its natural boundary. Here natural boundary means that no boundary
point has a neighborhood carrying a holomorphic function agreeing with
the indicated continuation on its domain side. Both statements concern
the entire lattice and all particle indices. They apply separately to
$\mathcal X=\beta^{-1}\chi$ and to $\chi$ with the field convention below.
\end{theorem}

The proof for unequal couplings is developed in the following sections.
Section~\ref{sec:closure} treats the isotropic slice by an exact reuse
interface, not by taking a limit in estimates whose constants degenerate
as $\lambda\to1$. No simultaneous uniformity in $\lambda$ is needed:
all local constants may depend on the fixed ratio and selected boundary
point. Rational and irrational ratios are included.

\subsection{Field normalization and state}
Let the magnetic field have energy units, so the finite-volume Hamiltonian
contains $-h\sum_v\sigma_v$, with $\sigma_v\in\{-1,1\}$ and magnetic
moment set to one. Direct differentiation of the finite partition
function gives
\[
 \left.\frac{\partial\langle\sigma_0\rangle}{\partial h}\right|_{h=0}
 =\beta\sum_v\bigl(\langle\sigma_0\sigma_v\rangle
                         -\langle\sigma_0\rangle\langle\sigma_v\rangle\bigr).
\]
Thus the normalized response is the covariance sum $\mathcal X$, and
$\chi=\beta\mathcal X$. A dimensionless field $\eta=\beta h$ instead
has response $\mathcal X$. A specified nonunit magnetic moment multiplies
both by its square.

On $0<z<z_c$ the Gibbs state is unique and its magnetization is zero:
\[
 \XH(z)=\sum_{(m,l)\in\Z^2}\langle\sigma_{00}\sigma_{ml}\rangle.
\]
On $z>z_c$ take the thermodynamic plus state, or equivalently the minus
state for this connected observable:
\[
 \XL(z)=\sum_{(m,l)\in\Z^2}
       \bigl(\langle\sigma_{00}\sigma_{ml}\rangle_+-M^2\bigr),
 \qquad M^2=(1-p^{-2})^{1/4}.
\]
The fourth root is positive on the displayed real interval.
For the energy-field response in this phase, the derivative is the
zero-field pure-phase linear response, equivalently the limit from
$h>0$ in the plus phase (or from $h<0$ in the minus phase).
The covariance formula is obtained with that phase fixed before taking
the zero-field limit. A symmetric
mixture has the same two-spin expectation but zero one-spin mean; using
its zero mean as the subtraction would leave the nondecaying term $M^2$
and would define a different, divergent sum. This distinction is retained
through the thermodynamic form-factor limit below.

The covariance sum also equals the literal uniform-field derivative after
the thermodynamic limit. Here is a precise justification once its finite
absolute sum, proved below, is available. The GHS inequality for arbitrary
nonnegative couplings and site fields \cite[Theorem 1.5]{CGN} says that
truncated two-spin correlations decrease as any nonnegative site field
increases. Pinning the spins removed from a larger plus box by fields
tending to $+\infty$ consequently shows that the plus-box covariances
increase with the box. Their sums converge monotonically to $\mathcal X$.
If $m_V(h)=\langle\sigma_0\rangle_{V,+,h}$, GHS also makes $m_V$
concave for $h\ge0$. Thus
\[
 m_V(h)-m_V(0)\le h m'_V(0)
 =\beta h\sum_{v\in V}\operatorname{Cov}_{V,+,0}(\sigma_0,\sigma_v)
 \le\beta h\mathcal X.
\]
Passing to the plus-state limit gives the upper right-derivative bound
$m'(0+)\le\beta\mathcal X$. Conversely, turn on $h$ only in a finite
set $E$ in the infinite plus state. This is a bounded local change of
measure, so its derivative at zero is exactly
$\beta\sum_{v\in E}\operatorname{Cov}_{+,0}(\sigma_0,\sigma_v)$.
Ferromagnetic field monotonicity makes the uniform-field magnetization
at least this local-field magnetization, with equality at $h=0$.
Letting $E$ increase to the lattice gives the reverse derivative bound.
Hence $m'(0+)=\beta\mathcal X$. In HT, spin-flip symmetry gives the
same derivative from the other side; in LT the stated one-sided
pure-phase convention applies. This argument does not exchange an
uncontrolled magnetic-field limit with an infinite sum.

\subsection{Published real-temperature form-factor interface}
Away from criticality on either positive real interval set
\[
 B=T/S>0,\qquad C=A/S,\qquad
 \cosh\gamma(q)=W(q)=C-B\cos q,\qquad \gamma(q)>0.
\]
Let $f_N(m,l)$ be
\begin{equation}\label{eq:corrFF}
 \frac{B^{N^2/2}}{N!}\int_{\T^N}
 \prod_{i=1}^N\frac{e^{-|m|\gamma(q_i)+\ii lq_i}}{\sinh\gamma(q_i)}
 \prod_{i<j}\left[
 \frac{\sin((q_i-q_j)/2)}{\sinh((\gamma(q_i)+\gamma(q_j))/2)}
 \right]^2\prod_{i=1}^N\frac{\dd q_i}{2\pi}.
\end{equation}
The exact infinite-plane correlation interface used here is
\begin{align}
 \langle\sigma_{00}\sigma_{ml}\rangle
   &=\xi_H\sum_{N\ge1,\ N\text{ odd}}f_N(m,l),
   &\xi_H&=\frac{(1-p^2)^{1/4}}{\sqrt p},\label{eq:HTcorr}\\
 \langle\sigma_{00}\sigma_{ml}\rangle_+
   &=\xi_L\left(1+\sum_{N\ge2,\ N\text{ even}}f_N(m,l)\right),
   &\xi_L&=(1-p^{-2})^{1/4}=M^2.\label{eq:LTcorr}
\end{align}
All powers are fixed by their positive real values. These are the
anisotropic vacuum-to-state spin form factors of
\cite{BL,IL}, in the bulk normalization of \cite{ONGP}. In particular,
$B^{N^2/2}$ is part of the $N$th term. Replacing it by a fixed prefactor
would change both the observable and the infinite-index problem.

For clarity about the thermodynamic use of these sources, the finite
cylinder formulas in \cite[(35)--(39)]{BL} concern arbitrary separations,
not just a single lattice row. Vacuum-to-state matrix elements are also
covered by \cite[Theorem 6.1]{IL} in the LT phase and by its
Remark 6.2, equations (6.8)--(6.10), with the HT vacuum change.
On a fixed real noncritical compact set, the one-body dispersion is
analytic in a common angular strip and has a positive real gap. The
finite-circumference correction functions $\eta$, $\Lambda^{-1}$ and
$\log\xi_T$ tend uniformly to zero: their sum-minus-integral expressions
are controlled by the exponentially convergent trapezoidal rule for
periodic analytic functions. They are uniformly bounded before the limit.

The particle limit is uniform as well. For real roots
$0<\rho(q)=e^{-\gamma(q)}\le r<1$, the Schur factor
$|(\rho_i-\rho_j)/(1-\rho_i\rho_j)|$ is at most a constant strictly below
one. The matched-Pfaffian argument proved in
Section~\ref{sec:regular} applies to the discrete angular grids as an
algebraic identity, with their normalized counting measures. Its
single-body finite-size corrections cost only $C^N$. Consequently the
finite-circumference coefficients, at each fixed separation, have a
common bound $\exp(-cN^2+O(N\log(N+1)))$, including their factorial
normalization. This is summable independently of circumference, so the
particle sum passes through the cylinder limit; taking each fixed-$N$
Riemann-sum limit gives \eqref{eq:corrFF}.

Every local cylinder limit is a translation-invariant, spin-flip invariant
Gibbs state for these finite-range positive couplings. Apply
\cite[Theorem 1]{Raoufi} with the vertex-transitive subgroup consisting
of translations: unequal horizontal and vertical couplings satisfy all
four stated hypotheses. Such a state is the equal mixture of the plus
and minus states. Its even spin correlations are those of the plus
state. This identifies the two-point function in \eqref{eq:LTcorr};
the connected pure-phase observable is obtained only after removing
the zero-particle term $\xi_L$.

\subsection{Symmetric double contours and every lattice site}
For separate contour variables $x,y$ put
\begin{equation}\label{eq:Dpoly}
 D(x,y;z)=A-\frac S2(x+x^{-1})-\frac T2(y+y^{-1}),\qquad
 \cA(v,w)=\frac{v-w}{1-vw}.
\end{equation}
With $X=\prod_i x_i$ and $Y=\prod_i y_i$, the same correlation
coefficient is
\begin{equation}\label{eq:doublecorr}
 \frac{p^{N/2}}{N!}\oint\cdots\oint
 X^{|m|-1}Y^{|l|-1}
 \prod_{i<j}\cA(x_i,x_j)\cA(y_i,y_j)
 \prod_i\frac{\dd x_i\,\dd y_i}{(2\pi\ii)^2D(x_i,y_i;z)}.
\end{equation}
The circles are positively oriented and sufficiently close to, but inside,
the unit circle. To check the normalization, fix $y=e^{\ii q}$ and
take the $x$ residue at the interior root $\rho=e^{-\gamma}$.
Since
\[
 D_x(\rho,y)=\frac{S(1-\rho^2)}{2\rho^2},\qquad
 \frac{\sqrt p}{D_x(\rho,y)}=\sqrt B\frac{\rho}{\sinh\gamma},
\]
and subtraction of two dispersion equations gives
\begin{equation}\label{eq:cancelpair}
 \cA(\rho_i,\rho_j)\cA(y_i,y_j)
 =-B\frac{(y_i-y_j)^2\rho_i\rho_j}
           {y_iy_j(1-\rho_i\rho_j)^2}
 =B\left[\frac{\sin((q_i-q_j)/2)}
               {\sinh((\gamma_i+\gamma_j)/2)}\right]^2=:P_{ij},
\end{equation}
the residue factors supply $B^{N/2}$ and the pair factors supply
$B^{N(N-1)/2}$. This recovers precisely $B^{N^2/2}$.
The remaining factors in \eqref{eq:doublecorr} give the stated Fourier
and spatial decay factors. Reflection in each coordinate is legitimate;
equality under interchange of the two separations at fixed unequal
couplings is not assumed.

The exact full-site multiplier in the double contour is
\begin{equation}\label{eq:allsites}
 \frac{(1+X)(1+Y)}{XY(1-X)(1-Y)}
 =\frac1{XY}+\frac{2(X^{-1}+Y^{-1})}{(1-X)(1-Y)}.
\end{equation}
It is the product of the two sums $\sum_{m\in\Z}X^{|m|-1}$ and
$\sum_{l\in\Z}Y^{|l|-1}$. It counts the origin once, each reflected
axis correctly, and all four open quadrants. If $C_N$ is the full-site
coefficient and $T_N$ the contour with the second numerator but without
its factor two, then $C_N=f_N(0,0)+2T_N$. The onsite sums are one in
the HT phase and $1-M^2$ after the LT subtraction. No axis or onsite
contribution is discarded.

Summing the Fourier variable and the geometric spatial series also gives
the complete angular representation
\begin{equation}\label{eq:fullangular}
 C_N=\frac1{N!}\int_{\sum q_i=0\bmod2\pi}
 G_N(q)\frac{1+Z}{1-Z}\prod_{i=1}^{N-1}\frac{\dd q_i}{2\pi},\quad
 G_N=B^{N/2}\prod_i\frac1{\sinh\gamma_i}\prod_{i<j}P_{ij},\quad
 Z=\prod_i e^{-\gamma_i}.
\end{equation}
There are exactly $N-1$ independent normalized angular measures, with
$q_N=-\sum_{i<N}q_i$ modulo $2\pi$.

\subsection{A holomorphic full sum on an independent spectral domain}
Define
\begin{equation}\label{eq:physicalspectralset}
 \cS_\lambda=\{z:\exists u,v\in[-1,1],\ A(z)=S(z)u+T(z)v\}.
\end{equation}
It is exactly the set where \eqref{eq:Dpoly} has a zero on the product
of two unit circles. Section~\ref{sec:geometry} proves that the components
of $\C\setminus\cS_\lambda$ containing $(0,z_c)$ and $(z_c,\infty)$
are exactly \eqref{eq:domains}, and that both are simply connected.

Here is the needed convergence argument independently of any proposed
natural-boundary conclusion. On a compact subset of the spectral
complement the torus polynomial has a nonzero minimum modulus. Move
both radii a small common amount inside one, within a product annulus
where this gap persists. The resulting contours can be held fixed on
a parameter neighborhood. Their Schur entries are bounded. The Schur
Pfaffian identity, followed by Hadamard's determinant inequality, bounds
the product of the two Pfaffians by
$C^N N^{(N+1)/2}$, with the harmless extra odd dimension obtained by
adjoining a constant ghost entry. After the $1/N!$ normalization and
the bounded one-body and full-site factors this is summable in $N$.
It is locally normally summable together with every fixed parameter
derivative by Cauchy's inequality on a slightly larger neighborhood.

On a physical real compact set the same contours in
\eqref{eq:doublecorr} additionally give
$r_x^{N|m|}r_y^{N|l|}$ with $r_x,r_y<1$. This proves absolute summability
in the lattice sites and particles and justifies their interchange.
Thus \eqref{eq:fullangular} and the double contour define the same
physical germs, with a proven full-site sum.

Write $F_N=p^{-N/2}C_N$ for the unscaled double contour, whose integrand
uses \eqref{eq:Dpoly} and \eqref{eq:allsites}. The complete sums are
\begin{align}
 \XH&=(1-p^2)^{1/4}\sum_{N\ge1,\ N\text{ odd}}p^{(N-1)/2}F_N,
                                                        \label{eq:globalHT}\\
 \XL&=(1-p^{-2})^{1/4}\sum_{N\ge2,\ N\text{ even}}p^{N/2}F_N.
                                                        \label{eq:globalLT}
\end{align}
The powers inside both sums are integers. The identity
$A^2=(1+S^2)(1+T^2)$ implies that every $p^2=1$ point lies in
$\cS_\lambda$. Also, $p$ never vanishes in $\OL\subset\{\Re z>0\}$;
in $\OH$ its only zero is the double zero at the origin. Hence the
prefactors in these formulas have unique holomorphic branches on the
specified simply connected domains, fixed by the real physical values.
Formula \eqref{eq:globalHT} removes the apparent HT $1/\sqrt p$
singularity at zero; $\XH(0)=1$. In particular, the one-particle term is
\[
 \XH^{(1)}=\frac{(1-p^2)^{1/4}}{A-S-T},
\]
which gives $1+2\beta(J_x+J_y)+O(\beta^2)$ at the origin.

For the angular formulas the spectral root is globally specified by
\[
 \rho=W-\sqrt{W^2-1},\qquad |\rho|<1,
 \qquad W=(A-T\cos q)/S.
\]
The square root has asymptotic $W$ at infinity. For $z$ in the spectral
complement, $W$ avoids $[-1,1]$ for every real $q$, so the inside root
is unique. At $S=0$ its root equation
$S\rho^2-2(A-T\cos q)\rho+S=0$ gives its removable value zero.
A local logarithm $\gamma=-\log\rho$ is used only where nonzero and
needed. No unsupported global logarithm or change of transfer vacuum
is imposed by periodic complex coupling shifts.

% END embedded section: sections/01_scope_physical.tex


% BEGIN embedded section: sections/02_geometry.tex
\section{Exact geometry on every fixed positive coupling slice}\label{sec:geometry}

Fix \(\lambda>0\), and use the one-complex-variable coordinate
\begin{gather*}
 z=2\beta J_x=x+\ii y,\qquad
 a(z)=\sinh z,\quad b(z)=\sinh(\lambda z),\\
 A(z)=\cosh z\cosh(\lambda z),\qquad p(z)=a(z)b(z).
\end{gather*}
The constraint \(K_y=\lambda K_x\) is imposed throughout. In particular,
the ratio \(b/a\) is generally a nonconstant meromorphic function of
\(z\); it must not be identified with the fixed coupling ratio
\(\lambda\).

\subsection{An independently specified spectral set}

\begin{definition}\label{def:spectralset}
For \((u,v)\in[-1,1]^2\), let
\[
 f_\lambda(z;u,v)=A(z)-a(z)u-b(z)v.
\]
The spectral set is
\begin{equation}\label{eq:spectralset}
 \cS_\lambda
 =\{z\in\C:\ f_\lambda(z;u,v)=0
                    \text{ for some }(u,v)\in[-1,1]^2\}.
\end{equation}
Equivalently, it is the set where the elementary spectral polynomial
\[
 A(z)-\frac{a(z)}2(X+X^{-1})
       -\frac{b(z)}2(Y+Y^{-1})
\]
has a zero with \(|X|=|Y|=1\).
\end{definition}

This definition uses neither the susceptibility nor a domain declared
maximal for that function. The anisotropic Nickel candidate equation
in \cite[eq.~(54)]{BHM} is the restriction of
\eqref{eq:spectralset} to cosine grids. A candidate equation, the singular
locus of a differential system, and a nonremovable singularity of a
particular infinite physical sum are distinct notions. Only the
geometric and grid-closure assertions are proved in this section.

Put
\begin{equation}\label{eq:cramerUV}
 D(z)=\Im(\overline{a(z)}b(z)),\qquad
 U(z)=\frac{\Im(\overline{A(z)}b(z))}{D(z)},\qquad
 V(z)=\frac{\Im(\overline{a(z)}A(z))}{D(z)}.
\end{equation}
When \(D\ne0\), the last two quantities are the unique real solution
of \(A=aU+bV\). Thus membership means \(|U|,|V|\le1\), and strict
inequalities give an interior point of \(\cS_\lambda\). At \(D=0\)
and \(a\ne0\), membership instead means
\[
 b/a\in\R,\qquad A/a\in\R,\qquad |A/a|\le1+|b/a|.
\]
If \(a=0\), the analogous condition with \(b\) applies. Simultaneous
zeros of \(a,b\) never belong to \(\cS_\lambda\), since then
\(A=\pm1\). These conditions retain the unsquared equation and its
sign information.

\begin{lemma}[Closure of both parity grid families]\label{lem:gridclosure}
For each parity separately, the closure of the union of zeros
\[
 f_\lambda\bigl(z;\cos(2\pi j/n),\cos(2\pi k/n)\bigr)=0
\]
is exactly \(\cS_\lambda\). Odd prime orders \(n=q\), and even orders
\(n=2q\) with \(q\) an odd prime, already suffice. The assertion remains
true after deleting grid pairs in any convention that leaves the grids
dense in \([-1,1]^2\), including deletion of the isotropic exceptional
pairs \(j=k=0\), and \(j+k=n/2\) when \(n\) is even.
\end{lemma}
\begin{proof}
The set \(\cS_\lambda\) is closed: for a convergent sequence of its
points, extract a convergent subsequence of representing pairs from the
compact square and pass to the limit in \(f_\lambda\).
This proves one inclusion.

For the reverse inclusion, fix \(z_0\in\cS_\lambda\) with representing
pair \((u,v)\), and choose allowed grid pairs \((u_m,v_m)\to(u,v)\).
The functions \(f_\lambda(\cdot;u_m,v_m)\) converge uniformly on compact
sets to \(f_\lambda(\cdot;u,v)\). The latter is not identically zero,
since its value at \(z=0\) is one. On a sufficiently small circle around
\(z_0\) it has no zeros, and Rouche's theorem supplies a zero of every
sufficiently close approximating function inside that circle. A sequence
of shrinking circles gives zeros converging to \(z_0\). This argument
also covers multiple roots, endpoints, \(D=0\), and zeros of either
hyperbolic sine.

The indicated prime grids are dense because their angular mesh tends
to zero. Removing the displayed exceptional pairs does not affect
density: when a nearest pair is deleted, changing one index by one
step, or a bounded number of steps, avoids the exceptional equalities
while the mesh error still tends to zero.
\end{proof}

The exclusions just mentioned are the convention displayed with the
isotropic equation in \cite[eq.~(52)]{BHM}. Lemma~\ref{lem:gridclosure}
does not assert that this list is an exhaustive anisotropic singularity
classification. The physical odd/even order restrictions are supplied
by the corresponding exact correlation representations, and the true
singularity and infinite-sum arguments are separate.

\begin{lemma}\label{lem:noisolatedS}
The set \(\cS_\lambda\) has no isolated points.
\end{lemma}
\begin{proof}
At a point \(z_0\in\cS_\lambda\), the coefficients \(a(z_0),b(z_0)\)
do not both vanish. Perturb a representing pair into the open square
so that \(f_\lambda(z_0;u_m,v_m)\ne0\), which is possible because that
affine function of the pair is not constant. The preceding Rouche
argument produces distinct nearby roots tending to \(z_0\).
\end{proof}

\subsection{A closed real-analytic inequality, including every degeneracy}

Define
\begin{align}
 F_1(x,y)&=\cosh^2(\lambda x)\cos^2y
                -\cosh^2x\sinh^2(\lambda x),\label{eq:F1}\\
 F_2(x,y)&=\cosh^2x\cos^2(\lambda y)
                -\sinh^2x\cosh^2(\lambda x).\label{eq:F2}
\end{align}

\begin{theorem}[Exact fixed-slice geometry]\label{thm:F1F2}
For every \(\lambda>0\),
\begin{equation}\label{eq:F1F2set}
 \boxed{\cS_\lambda=\{x+\ii y:\ F_1(x,y)F_2(x,y)\le0\}.}
\end{equation}
\end{theorem}
\begin{proof}
Reflection in the imaginary axis preserves the set and the two real
functions, so it is enough to consider \(x\ge0\). At \(x=0\),
\(a,b\) are purely imaginary and \(A\) is real. Membership is exactly
\[
 A=\cos y\cos(\lambda y)=0,
\]
which is precisely \(F_1F_2=0\) there.

Suppose now \(x>0\), and abbreviate
\begin{gather*}
 h=\sinh x,\quad c=\cosh x,\quad H=\sinh(\lambda x),\quad C=\cosh(\lambda x),\\
 t=\cos y,\quad q=\sin y,\quad T=\cos(\lambda y),\quad Q=\sin(\lambda y).
\end{gather*}
Write \(N_U=\Im(\overline A b)\), \(N_V=\Im(\overline a A)\).
Expansion gives the real identities
\[
 D=hCtQ-cHqT,\qquad
 N_U=ctTQ-hHCq,\qquad N_V=hcHQ-CtqT,
\]
and the exact factorizations
\begin{align}
 N_U^2-D^2&=(c^2T^2-h^2C^2)(t^2Q^2-H^2q^2),\label{eq:NUfactor}\\
 N_V^2-D^2&=(c^2H^2-C^2t^2)(h^2Q^2-q^2T^2).\label{eq:NVfactor}
\end{align}
For clarity set
\[
 a_0=h^2,\ b_0=H^2,\ A_0=1+a_0,\ B_0=1+b_0,
\quad X=t^2,\ Y=T^2,\quad
 X_0=\frac{A_0b_0}{B_0},\quad Y_0=\frac{a_0B_0}{A_0}.
\]
Then \(F_1=B_0(X-X_0)\), \(F_2=A_0(Y-Y_0)\), and
\begin{align}
 N_U^2-D^2&=F_2 K,& K&=(B_0-Y)X-b_0,\label{eq:KF}\\
 N_V^2-D^2&=-F_1 L,& L&=a_0-(A_0-X)Y.\label{eq:LF}
\end{align}
Both \(K,L\) vanish at \((X_0,Y_0)\), increase with \(X\), and
decrease with \(Y\). Thus \(F_1>0,F_2<0\) implies \(K,L>0\),
whereas \(F_1<0,F_2>0\) implies \(K,L<0\). In either case the
right-hand sides in \eqref{eq:KF}--\eqref{eq:LF} are negative.
In particular \(D\ne0\), since otherwise a real square would be
negative, and \(|U|,|V|<1\). This proves membership for strict
opposite signs without an illegitimate division by \(D\).

If \(D\ne0\), \(0<Y<1\), and both \(F\)'s are positive, the
two membership inequalities would require
\[
 X\le\frac{b_0}{B_0-Y},\qquad X\ge A_0-\frac{a_0}{Y}.
\]
They are inconsistent because \(Y>Y_0\) and
\begin{equation}\label{eq:rationaldifference}
 \frac{b_0}{B_0-Y}-\left(A_0-\frac{a_0}{Y}\right)
 =\frac{A_0(Y-1)(Y-Y_0)}{Y(B_0-Y)}<0.
\end{equation}
If both \(F\)'s are negative, the required inequalities reverse;
now \(Y<Y_0\) makes \eqref{eq:rationaldifference} positive, again
inconsistent. At \(Y=0\), the factors directly give membership
exactly when \(X\ge X_0\). At \(Y=1\), they give membership exactly
when \(Y_0\le1\) and \(X\le X_0\). These conditions agree with
\(F_1F_2\le0\); if \(Y_0>1\), then also \(X_0>1\), so the opposite
sign possibility at \(Y=1\) cannot occur.

It remains to handle zeros and collinearity. Every point with \(F_1=0\)
belongs to the spectral set, including those with \(D=0\). Indeed, the
two endpoint dispersion root pairs are
\begin{align}
 r_+(z)&=e^{\lambda z}\tanh(z/2),&
 \frac{r_++r_+^{-1}}2&=\frac{A-b}{a},\label{eq:endpointplus}\\
 r_-(z)&=e^{-\lambda z}\tanh(z/2),&
 \frac{r_-+r_-^{-1}}2&=\frac{A+b}{a}.\label{eq:endpointminus}
\end{align}
Here \(a\ne0\), and \(\tanh(z/2)\) has neither zero nor pole because
\(x>0\). The exact modulus identity
\[
 |\tanh(z/2)|^2=\frac{\cosh x-\cos y}{\cosh x+\cos y}
\]
shows that \(|r_+|=1\) exactly when
\(\cos y=\cosh x\tanh(\lambda x)\), and \(|r_-|=1\) exactly when
the right side has a minus sign. These are the two possibilities in
\(F_1=0\). A unit-modulus root has \((r+r^{-1})/2\in[-1,1]\),
giving respectively the representing endpoint \(V=1\) or \(V=-1\).
Exchanging the two axes proves the same assertion for \(F_2=0\).

Conversely suppose \(D=0\) at a candidate point. Then \(N_U=N_V=0\).
If both \(F\)'s were nonzero, \eqref{eq:KF}--\eqref{eq:LF} would give
\(K=L=0\). These force \(Y>0\), and
\eqref{eq:rationaldifference} forces \(Y=Y_0\) or \(Y=1\).
The first contradicts \(F_2\ne0\). In the second case \(X=1\) too.
The coefficients are then real up to their signs, and candidate
membership requires
\[
 \cosh x\cosh(\lambda x)\le\sinh x+\sinh(\lambda x).
\]
The square of the left side minus the square of the right side is
\((\sinh x\sinh(\lambda x)-1)^2\). Equality is therefore forced,
and \(F_1=F_2=0\), again a contradiction. Finally, at \(F_1=0\)
identity \eqref{eq:KF} becomes
\(N_U^2-D^2=-(b_0/B_0)F_2^2\); at \(F_2=0\),
\eqref{eq:LF} becomes \(N_V^2-D^2=-(a_0/A_0)F_1^2\).
Thus a collinear candidate with one zero has both zeros. All cases are
covered, proving \eqref{eq:F1F2set}.
\end{proof}

\subsection{The two exact physical spectral components}

Let \(z_c>0\) be the unique solution
\begin{equation}\label{eq:zc}
 \sinh z_c\sinh(\lambda z_c)=1.
\end{equation}
Define strictly increasing functions on \([0,\infty)\) by
\begin{equation}\label{eq:PQ}
 P_\lambda(t)=\cosh t\tanh(\lambda t),\qquad
 Q_\lambda(t)=\cosh(\lambda t)\tanh t.
\end{equation}
Both start at zero, tend to infinity, and satisfy
\(P_\lambda(z_c)=Q_\lambda(z_c)=1\). Their real inverse functions are
unambiguous. Put
\begin{align}
 Y_\lambda&=\frac{\pi}{2\max(1,\lambda)},\label{eq:Ylambda}\\
 H_\lambda(y)&=\min\{P_\lambda^{-1}(\cos y),
                       Q_\lambda^{-1}(\cos(\lambda y))\},
                       &&|y|<Y_\lambda,\label{eq:Hlambda}\\
 L_\lambda(y)&=\max\{P_\lambda^{-1}(|\cos y|),
                       Q_\lambda^{-1}(|\cos(\lambda y)|)\},
                       &&y\in\R.\label{eq:Llambda}
\end{align}

\begin{theorem}[Physical spectral components]\label{thm:components}
The component of \(\C\setminus\cS_\lambda\) containing the real
high-temperature interval \((0,z_c)\), and the component containing
the real low-temperature interval \((z_c,\infty)\), are respectively
\begin{align}
 \OH&=\{x+\ii y:\ |y|<Y_\lambda,
                         \ |x|<H_\lambda(y)\},\label{eq:OH}\\
 \OL&=\{x+\ii y:\ x>L_\lambda(y)\}.\label{eq:OL}
\end{align}
Both are simply connected. Their boundaries are the complete lens
boundary and the complete graph \(x=L_\lambda(y)\), respectively.
Moreover \(0\le L_\lambda(y)\le z_c\), and
\(\cS_\lambda\subset\{|\Re z|\le z_c\}\).
\end{theorem}
\begin{proof}
At the origin \(F_1,F_2\) are both positive. Their signs cannot change
in its complement component. Since neither cosine can then vanish,
both retain their positive signs, giving
\(|y|<Y_\lambda\). The two positive inequalities are exactly
\(|x|<H_\lambda(y)\). Conversely this entire lens is connected:
decrease \(|x|\) to zero at fixed \(y\), then decrease \(|y|\) to zero.
Both operations preserve the inequalities. In fact the lens is star
shaped about zero and hence simply connected.

At real \(z>z_c\), both \(F\)'s are negative. For \(x>0\), their
simultaneous negativity is exactly \(x>L_\lambda(y)\). This epigraph
is connected because \(L_\lambda\le z_c\), and a component with both
signs negative cannot cross \(x=0\), where both are nonnegative.
Thus it is precisely the LT component. Its continuous boundary graph
is mapped to a straight line by the homeomorphism
\((x,y)\mapsto(x-L_\lambda(y),y)\), proving simple connectivity.
Finally if \(|x|>z_c\), both \(P_\lambda(|x|),Q_\lambda(|x|)\)
exceed one, so both \(F\)'s are negative and the point is outside
\(\cS_\lambda\).
\end{proof}

The boundary functions are explicitly defined before any analyticity
or singularity assertion for a particle sum. In particular the set
\(\cS_\lambda\) itself generally has interior, whereas the boundaries
in Theorem~\ref{thm:components} are one-dimensional barriers.

\begin{proposition}[Strict anisotropic opening]\label{prop:anisotropicarea}
For every \(\lambda\ne1\), \(\cS_\lambda\) has nonempty interior
arbitrarily close to \(z_c\). The horizontal width between the two
local endpoint curves near \(z_c+\ii y\) is
\begin{equation}\label{eq:cuspwidth}
 \frac12\bigl|\lambda\coth(\lambda z_c)-\coth z_c\bigr|y^2
       +O(y^4).
\end{equation}
\end{proposition}
\begin{proof}
For small \(y\), the two zero curves are
\[
 x_1(y)=P_\lambda^{-1}(\cos y)
      =z_c-\frac{y^2}{2P_\lambda'(z_c)}+O(y^4),
\]
\[
 x_2(y)=Q_\lambda^{-1}(\cos(\lambda y))
      =z_c-\frac{\lambda^2y^2}{2Q_\lambda'(z_c)}+O(y^4).
\]
Using \(\sinh z_c\sinh(\lambda z_c)=1\), subtraction of the two
displayed coefficients gives \eqref{eq:cuspwidth}. The coefficient
is nonzero for \(\lambda\ne1\): the function
\(t\mapsto t\coth t\) is strictly increasing for \(t>0\), since
\(\sinh t\cosh t-t>0\). Strictly between the two distinct curves
the \(F\)'s have opposite signs, so Theorem~\ref{thm:F1F2} gives
an open part of the spectral set.
\end{proof}

For a concrete exact check, at \(\lambda=2\) take
\[
 a_0=\frac{1+\ii\sqrt7}{4},\qquad c_0=\sqrt{1+a_0^2},\quad
 \Re c_0>0,\qquad z_0=\operatorname{arsinh}a_0
\]
on the corresponding local branch. Then \(b=2a_0c_0\),
\(A=c_0(1+2a_0^2)=b/2\), so \(U=0,V=1/2\).
Also \(D=\Im c_0\ne0\) and \(|ab|=2^{-1/4}\ne1\).
Thus the frequently suggested replacement of the entire candidate
set by \(|ab|=1\) is already false geometrically. This observation is
not a refutation of any physical natural-boundary theorem.

\subsection{Zeros, collinearity and every exceptional boundary point}

The imaginary-axis intersection has the exact simple description
\begin{equation}\label{eq:imagaxisS}
 \cS_\lambda\cap\ii\R
 =\{\ii y:\cos y\cos(\lambda y)=0\}.
\end{equation}
It is locally finite. Both cosines vanish simultaneously if and only
if \(\lambda\) is a ratio of two positive odd integers. In reduced
form \(\lambda=r/s\) with both \(r,s\) odd, these simultaneous
points have \(y=(2m+1)\pi s/2\).

For completeness, a zero of \(a\) has \(z=\ii\pi m\) and belongs
to \(\cS_\lambda\) exactly when \(\cos(\lambda\pi m)=0\);
then \(b=\pm\ii\), \(A=0\), and its representing pairs have
\(V=0\). The corresponding statement for \(b=0\) follows by
axis exchange. Simultaneous sine zeros are outside the spectral set.
Inside \(\OH\), the only sine zero is the common zero at \(z=0\);
inside \(\OL\) neither sine vanishes. In addition,
\begin{equation}\label{eq:cornerp}
 \bigcup_{\eps_1,\eps_2\in\{1,-1\}}
 \{A=\eps_1 a+\eps_2b\}=\{p^2=1\}\subset\cS_\lambda,
\end{equation}
because
\[
 A^2-(\eps_1 a+\eps_2 b)^2=(p-\eps_1\eps_2)^2.
\]
Every such set is discrete on the fixed nonconstant slice. Hence
the phase prefactors have no zero or branch point in either physical
component, except for the removable HT normalization issue at the
origin treated by its integer-power formula. On \(\OL\), even that
issue is absent. The boundary of \(\OL\) reaches the imaginary axis
only at simultaneous cosine zeros.

\begin{lemma}[Collinear candidates]\label{lem:collinear}
For \(\Re z\ne0\),
\[
 \{z\in\cS_\lambda:D(z)=0\}
 =\{z:D(z)=0,\ |a(z)b(z)|=1\}.
\]
For each \(\lambda>0\), \(\lambda\ne1\), this set, together with
the collinear candidates on the imaginary axis, is locally discrete.
\end{lemma}
\begin{proof}
Off the imaginary axis, \(a,b\ne0\). Write \(b=r a\), \(r\in\R\setminus
\{0\}\). A candidate also has \(A=t a\), \(t\in\R\), and
\[
 t^2=1+r^2+r^2a^2+a^{-2}.
\]
If \(a=R e^{\ii\theta}\), reality of this expression implies
\((r^2R^2-R^{-2})\sin(2\theta)=0\). Unless \(a\) is real or purely
imaginary, this gives \(|ab|=1\). If \(a,b\) are real, the inequality
\(|A|\le|a|+|b|\), together with
\[
 |A|^2-(|a|+|b|)^2=(|ab|-1)^2,
\]
gives the same conclusion. If both are purely imaginary off the
imaginary axis, then both cosines vanish, \(|a|,|b|>1\), and \(A\)
is nonzero real; it cannot be a real multiple of imaginary \(a\).
Conversely, if \(D=0\) and \(|ab|=1\), then
\[
 (A/a)^2=1+r^2+2|r|\cos(2\theta)
                   \in[(1-|r|)^2,(1+|r|)^2].
\]
Thus \(A/a\) is real with the required segment bound, proving the
set equality.

By symmetry it is enough to prove local discreteness for \(x>0\).
For each \(y\) with \(|\sin y\sin(\lambda y)|<1\), the equation
\begin{equation}\label{eq:modulusgraph}
 (\sinh^2x+\sin^2y)(\sinh^2(\lambda x)+\sin^2(\lambda y))=1
\end{equation}
has a unique solution \(x=X(y)>0\). The left side strictly increases
with \(x\), starts below one, and tends to infinity. Its positive
\(x\) derivative proves that \(X\) is real analytic. At a simultaneous
cosine zero it tends continuously to zero.

If there are no simultaneous cosine zeros, \(X\) is analytic on all
of \(\R\). The analytic function \(D(X(y)+\ii y)\) is not identically
zero, because its derivative at \(y=0\) is
\[
 a(z_c)b'(z_c)-a'(z_c)b(z_c)
       =\lambda\coth(\lambda z_c)-\coth z_c\ne0.
\]
Its zeros are therefore locally discrete. If simultaneous cosine
zeros occur, write \(\lambda=r/s\) in reduced odd/odd form. The
intervals between \((2m+1)\pi s/2\) are translates by \(\pi s\) of
one another, and each contains a translate of \(y=0\). Both \(X\)
and the displayed restriction of \(D\) have that period, so the
restriction is not identically zero on any interval.

There is no accumulation of these zeros at the interval endpoints.
At a simultaneous zero, write \(z=z_0+\zeta\), \(\zeta=x+\ii\delta\).
For signs \(\sigma,\tau\in\{1,-1\}\),
\[
 a=\ii\sigma\cosh\zeta,\qquad
 b=\ii\tau\cosh(\lambda\zeta),\qquad
 D=\sigma\tau(\lambda^2-1)x\delta+O(|\zeta|^4).
\]
Equation \eqref{eq:modulusgraph} gives
\(X(y)=|\delta|+O(|\delta|^3)\), so \(D\ne0\) for all sufficiently
small nonzero \(\delta\). At any other imaginary-axis candidate,
\(|ab|<1\), and continuity excludes approach by off-axis collinear
candidates with \(|ab|=1\). This proves local discreteness everywhere.
\end{proof}

\begin{definition}\label{def:regularedge}
A regular edge point has \(D\ne0\), exactly one of \(U,V\) equal to
\(1\) or \(-1\), the other strictly in \((-1,1)\), and nonzero
derivative of the corresponding meromorphic edge map
\[
 h_\eps(z)=\frac{A(z)-\eps a(z)}{b(z)}
 \quad(U=\eps),\qquad
 g_\eps(z)=\frac{A(z)-\eps b(z)}{a(z)}
 \quad(V=\eps).
\]
\end{definition}

\begin{proposition}[Dense regular boundary arcs]\label{prop:regular_dense}
For every fixed \(\lambda>0\), \(\lambda\ne1\), regular edge points
are dense in the boundary of every component of
\(\C\setminus\cS_\lambda\), in particular in both boundaries in
Theorem~\ref{thm:components}. The excluded set is locally discrete.
\end{proposition}
\begin{proof}
The four edge maps are nonconstant meromorphic functions: each has a
pole at \(z=0\). Their critical points and poles are locally discrete.
Corner points are the discrete zeros in \eqref{eq:cornerp}; collinear
candidates are locally discrete by Lemma~\ref{lem:collinear}. Let
\(E\) be their union, restricted to \(\cS_\lambda\).

At a boundary point outside \(E\), strict inequalities cannot hold
for both coordinates, since they would give an interior point.
Exactly one is an endpoint, and the corresponding holomorphic edge
map has nonzero derivative. Its imaginary part is a nondegenerate
local defining function for that edge. The unique real-coordinate
solution in \eqref{eq:cramerUV} then shows that \(\cS_\lambda\) is
locally one closed side of an analytic arc and its complement is
the other side.

If \(\Omega\) is any complement component, then
\(\partial\Omega\subset\partial\cS_\lambda\). It has no isolated
point: if \(z_0\) were isolated in \(\partial\Omega\), a sufficiently
small punctured disk, being connected and containing points of
\(\Omega\), would lie wholly in \(\Omega\). Then \(z_0\) would be
an isolated point of \(\cS_\lambda\), contrary to
Lemma~\ref{lem:noisolatedS}. Removing the locally discrete set \(E\)
from a boundary without isolated points leaves a dense subset.
\end{proof}

At the right side of the physical HT lens, both cosines in
\eqref{eq:Hlambda} are positive. Equations
\eqref{eq:endpointplus}--\eqref{eq:endpointminus}, or their axis-swapped
versions, show that every regular edge there is a positive endpoint.
Reflection \(z\mapsto-\overline z\) preserves the lens and sends the
representing pair \((U,V)\) to \((-U,-V)\), so every regular left edge
is a negative endpoint. Thus the negative-endpoint issue for the odd
sum occurs on the actual physical component for every
\(\lambda\ne1\); it cannot be dismissed as a remote-sheet issue.
On the LT boundary, an \(F_1=0\) regular edge has
\(V=\operatorname{sign}(\cos y)\), and an \(F_2=0\) regular edge has
\(U=\operatorname{sign}(\cos(\lambda y))\). These alternatives cover
all its regular pieces.

\subsection{The canonical root and the precise resonance arithmetic}

\begin{lemma}[Canonical spectral sheet]\label{lem:canonicalroot}
For \(z\notin\cS_\lambda\) and real \(q\), the equation
\[
 a\rho^2-2(A-b\cos q)\rho+a=0
\]
has a unique root with \(|\rho|<1\), interpreted as \(\rho=0\) if
\(a=0\). It is holomorphic in \(z\), locally jointly holomorphic
with the complex angle, and has a common modulus gap below one on
every compact subset of \(\C\setminus\cS_\lambda\) and the real
angle circle. On the positive real HT and LT intervals it is the
positive physical decaying root.
\end{lemma}
\begin{proof}
When \(a\ne0\), put \(W=(A-b\cos q)/a\). Outside the spectral set,
\(W\notin[-1,1]\). The root
\[
 \rho(W)=W-\sqrt{W^2-1},\qquad \sqrt{W^2-1}\sim W\quad(W\to\infty),
\]
is the unique interior root on \(\C\setminus[-1,1]\); the other root
is its reciprocal. At \(a=0\), the coefficient \(A-b\cos q\ne0\)
by the spectral-set definition, so the implicit-function theorem
extends the root through \(\rho=0\). Compactness supplies the uniform
gap and a thin complex angle neighborhood. For positive real
\(z\ne z_c\), \(A-a-b>0\), since
\[
 (A-a-b)(A+a+b)=(ab-1)^2,
\]
and hence \(W(q)>1\) for all real \(q\), proving the physical match.
\end{proof}

The root statement does not identify the two particle parities or
their phase prefactors. It also does not assert a logarithm at a root
zero. On the physical LT component all these zeros are absent and
simple connectivity fixes the nonvanishing algebraic prefactors by
their real physical values. On the HT component the only root-zero
normalization issue is at the origin, handled in the physical section.

At a regular edge with one cosine equal to \(+1\) and the other
equal to \(\cos(2\pi k/q)\), where \(q\) is an odd prime and
\(1\le k\le(q-1)/2\), uniqueness in \eqref{eq:cramerUV} implies that
the least grid order is \(q\) and the least even order is \(2q\).
These values are dense on every regular edge arc because its
nonramified edge coordinate maps the arc to a real interval. At a
\(-1\) endpoint, an exact grid representation requires an even
order, and the least even order is again \(2q\). There is no exact
odd endpoint resonance on such an edge. These arithmetic facts alone
do not prove lower-form-factor regularity or noncancellation; the
actual integral arguments provide those conclusions separately.

One useful exact accumulation from the forbidden side is the following.
At a regular \(U=-1\) point with
\(V=\cos(2\pi k/q)\), take odd \(n=q(2m+1)\) and set
\[
 u_n=-\cos(\pi/n),\qquad v_n=V.
\]
The simple root of the candidate equation has
\[
 z_n=z_0+\frac{a(z_0)}{\partial_z f_\lambda(z_0;-1,V)}
                  \frac{\pi^2}{2n^2}+O(n^{-4}).
\]
It lies in the interior of \(\cS_\lambda\) for large \(n\), and its
least grid order is \(n\), since \((n-1)/2\) is coprime to \(n\).
The derivative in the denominator is nonzero because the regular
edge map is nonramified. This is candidate accumulation only; no
continuation of the infinite sum into that interior is assumed.

\subsection{Isotropic covering and ratio symmetries}

At \(\lambda=1\), \(a=b=s=\sinh z\) and \(A=1+s^2\), so
\[
 A/a=s+s^{-1}=u+v\in[-2,2].
\]
The preimage of \([-2,2]\) under the Joukowski map \(s+s^{-1}\)
is exactly \(|s|=1\). Thus
\begin{equation}\label{eq:isotropicgeometry}
 \cS_1=\{z:|\sinh z|=1\}.
\end{equation}
Equivalently, \(F_1=F_2\), and their common zero is
\(\cos^2 y=\sinh^2x\). Since \(P_1=Q_1=\sinh\), the two physical
components become
\[
 \OH(1)=\{|y|<\pi/2,\ |x|<\operatorname{arsinh}(\cos y)\},
 \qquad
 \OL(1)=\{x>\operatorname{arsinh}(|\cos y|)\}.
\]
The map \(s=\sinh z\) maps the first biholomorphically onto the unit
disk, with the inverse-sinh branch real at the origin. It maps the
second as the universal cover of \(|s|>1\), with deck translations
\(z\mapsto z+2\pi\ii\). To see the latter directly, put \(w=e^z\).
For \(|s|>1\), the equation \(w-w^{-1}=2s\) has exactly one root
with \(|w|>1\), depending holomorphically on \(s\). A unit-modulus
root would instead give \(s\in\ii[-1,1]\). The exterior \(w\) domain
is a continuous radial epigraph; exponentiation from \(\OL(1)\)
is its universal cover. This proves the covering assertion.

At the boundary points with \(\cosh z=0\), the inverse-sinh coordinate
is branched. Elsewhere on the boundary it is locally biholomorphic.
Accordingly the inherited dense isotropic singular family can be
transferred at nonramified points and then closed by density; there
is no need to assume an inverse-sinh branch through a branch point.
This is the interpretation of \(|\sinh(2\beta J)|=1\) used here.

Finally \(\cS_\lambda\) is invariant under conjugation and negation,
and axis exchange gives
\[
 \cS_{1/\lambda}=\lambda\cS_\lambda.
\]
If \(\lambda=r/s\) is rational in lowest positive terms,
\(\ii\pi s\) is an imaginary period of the spectral set: the two
representing cosines only change signs. This statement does not
automatically give a period of the physical odd susceptibility.
For irrational \(\lambda\), there is no nonzero imaginary period.
Indeed \eqref{eq:imagaxisS} is the union of two arithmetic progressions
with spacings \(\pi\) and \(\pi/\lambda\). A translate of the first
progression can intersect the second in at most one point unless
their spacing ratio is rational. A period of the union must therefore
be a period of each progression, forcing it to belong to both
\(\pi\Z\) and \((\pi/\lambda)\Z\), hence to be zero.

The exceptional ratio \(\lambda=1\), the odd/odd rational simultaneous
cosine zeros, the remaining sine or cosine zeros, corner points and
meromorphic critical points have all been accounted for above.
For \(\lambda\ne1\), none removes a boundary interval from the dense
regular-arc conclusion. Which complete physical sums have genuine
singularities on those arcs is the task of the following sections.

% END embedded section: sections/02_geometry.tex


% BEGIN embedded section: sections/03_regular_edges.tex
\section{Regular spectral edges: whole amplitudes and the entire tail}
\label{sec:regular}

At a regular boundary point exactly one of the two real spectral
coordinates is $+1$ or $-1$, the other lies strictly between them, and
the corresponding edge map is a nonramified holomorphic local
coordinate. Interchange the transfer axes when necessary. First take
the endpoint $q_0=0$ and write
\begin{equation}\label{eq:positiveedge}
 W(q;u)=b+u+B(u)(1-\cos q),\qquad
 b=\cos\phi,\quad 0<\phi<\pi,\quad \Im B(0)>0.
\end{equation}
Here $u=W(0)-b$ vanishes at the selected point. The opposite sign of
$\Im B$ is obtained by complex conjugation. On the edge, $W(q)$ is
in the upper half-plane except at $q=0$. Its canonical root
$\rho(q)=e^{-\gamma(q)}$ lies in the lower half-disk, with
$\rho(0)=e^{-\ii\phi}$. Compactness gives
\begin{equation}\label{eq:halfdisc}
 \Im\rho(q)\le-\delta<0,\qquad |\rho(q)|\le1
 \quad(q\in\R).
\end{equation}
There is no real-angle point with $W=\pm1$. The root, a local
$\gamma$, and all one-body and canceled-pair denominators are therefore
holomorphic on a common fixed complex parameter and angular
neighborhood. On the physical-domain approach $u=\ii\eps$, $\eps>0$,
\begin{equation}\label{eq:attenuation}
 \Re\gamma(q;\ii\eps)\ge c\bigl(\eps+\dist(q,0)^2\bigr).
\end{equation}
Indeed $\Im W=\eps+\Im B(\ii\eps)(1-\cos q)$ is bounded below by
that expression, while
$\Im W=\sinh(\Re\gamma)\sin(\Im\gamma)$ on the compact chosen branch.
The same local side belongs to the physical component bordering the
regular edge, by the geometry of Section~\ref{sec:geometry}.

\subsection{A strict Schur bound, with the anisotropy retained}
The elementary identity
\begin{equation}\label{eq:schurdifference}
 |1-vw|^2-|v-w|^2
 =(1-|v|^2)(1-|w|^2)+4\Im v\Im w
\end{equation}
and \eqref{eq:halfdisc} give a constant $\rho_0<1$ such that
$|\cA(\rho_i,\rho_j)|\le\rho_0$ for every pair of real angles.
By continuity a slightly larger constant still below one works on a
fixed complex parameter neighborhood. This assertion continues to hold
there even where a separate inequality $|\rho_i|\le1$ no longer does.
Every complete canceled pair $P_{ij}$ is analytic and bounded on the
same neighborhood.

It is not true that $|P_{ij}|\le1$ there. For example, on
\eqref{eq:positiveedge},
\[
 P(q,-q)=\frac{B\sin^2q}
               {[b+B(1-\cos q)]^2-1}.
\]
When $\Im B$ is small and the real part of the bracket is close to one,
this can be arbitrarily large. The proof below retains the strict bound
only for the uncanceled $\rho$-Schur factors.

\begin{lemma}[Matched-Pfaffian numerator estimate]\label{lem:matching}
For every fixed integer $r\ge0$, on a common parameter disk,
\[
 |\partial_u^rG_N(q;u)|
 \le C_r\exp\{-\kappa N^2+C N\log(N+1)\}
\]
uniformly in the real angles and all positive $N$. The constants may
depend on the edge but not on $N$.
\end{lemma}
\begin{proof}
For even $N$, expand only the Schur Pfaffian in the $y_i=e^{\ii q_i}$
variables. Each of its $(N-1)!!$ matchings contains $N/2$ edges. Combine
each matched $y$ factor with its corresponding $\rho$ factor using
\eqref{eq:cancelpair}; bound that complete pair by a constant $C$.
The other $N(N-2)/2$ factors are bounded by $\rho_0$. Thus
\[
 \left|\prod_{i<j}P_{ij}\right|
 \le (N-1)!!\,C^{N/2}\rho_0^{N(N-2)/2}.
\]
For odd $N$, adjoin a ghost variable $-1$ to the $y$ Pfaffian.
Its Schur entry with every real variable is the removable constant one.
There are $N!!$ matchings, with $(N-1)/2$ real matched edges and one
ghost edge, leaving $(N-1)^2/2$ uncanceled $\rho$ factors. This gives
\[
 \left|\prod_{i<j}P_{ij}\right|
 \le N!!\,C^{(N-1)/2}\rho_0^{(N-1)^2/2}.
\]
The identities initially hold away from reciprocal $y$ pairs and extend
by the analytic canceled formula. No division by a vanishing $\rho$
factor occurs. One-body factors, including $B^{N/2}$, cost $C^N$.
These are Gaussian bounds in $N$ up to $\exp(O(N\log(N+1)))$.
Cauchy's inequality on the common disk proves the numerator-jet
statement. It is applied only to $G_N$, not across a pole of
$(1+Z)/(1-Z)$.
\end{proof}

\subsection{All indices in the differentiated tail}
Put $Q_0(q)=\sum_i\dist(q_i,0)^2$. Equation
\eqref{eq:attenuation} implies
\[
 |1-Z|\ge1-|Z|\ge c\min\{1,N\eps+Q_0(q)\}.
\]
All fixed parameter derivatives of $\gamma$ are bounded. Differentiating
the exponential and rational factor consequently gives, for every fixed
$j$ and $0\le r\le j$,
\begin{equation}\label{eq:globaljet}
 \left|\partial_u^r\frac{1+Z}{1-Z}\right|
 \le C_j N^j\bigl(1+Q_0(q)^{-j-1}\bigr).
\end{equation}
For $N-1\ge2j+3$, retain just the first $2j+3$ independent angle
squares in $Q_0$. Their inverse $(j+1)$st power is integrable in
$2j+3$ real dimensions, with a constant independent of $N$; all the
unused normalized angular measures have mass one. Lemma~\ref{lem:matching}
therefore implies
\begin{equation}\label{eq:hightail}
 |\partial_u^j C_N(\ii\eps)|
 \le \frac{C_jN^j}{N!}
       \exp\{-\kappa N^2+C N\log(N+1)\},\qquad N\ge2j+4.
\end{equation}
This sums uniformly in $\eps$.

For any remaining fixed $N$, the only real configuration at which
$|Z|=1$ is $q_1=\cdots=q_N=0$. Near it, set
$q_N=-\sum_{i<N}q_i$, $Q=\sum_iq_i^2$ and
$\Delta(q)=\prod_{i<j}(q_i-q_j)$. The numerator and every fixed
parameter jet have the factor $\Delta^2$. In radial coordinates the
worst $j$th derivative is bounded by an integrand proportional to
\[
 r^{N(N-1)}r^{N-2}(\eps+r^2)^{-j-1}\dd r.
\]
Its limiting exponent is $N^2-2j-4$. If $n$ is the first resonant
index of the chosen parity, take $k=(n^2-1)/2$ for odd $n$ or
$k=n^2/2-1$ for even $n$. For every larger same-parity $N$ and
$j\le k$ this exponent is greater than $-1$. Combining the finitely
many such indices below $2j+4$ with \eqref{eq:hightail} proves
\begin{equation}\label{eq:alltail}
 \sum_{\substack{N>n\\N\text{ of the selected parity}}}
 |\partial_u^j C_N(\ii\eps)|=O_j(1),\qquad0\le j\le k.
\end{equation}
This controls the entire differentiated tail, with no large-index sector
left implicit.

\subsection{The first amplitude of the complete form factor}
Choose $\phi=2\pi a/p_0$ with $p_0$ an odd prime and
$1\le a\le(p_0-1)/2$. Then the first odd resonance is $n=p_0$ and
the first even resonance is $n=2p_0$. For all lower indices of the
selected parity, the only possible unit-modulus product configuration
has $Z=e^{-\ii N\phi}\ne1$, so those entire form factors are
holomorphic on a neighborhood of the selected point.

The Gaussian Vandermonde period in the actual independent-angle measure
is
\begin{equation}\label{eq:mehta}
 G_n^{\rm Gauss}:=
 \int_{q_n=-\sum_{i<n}q_i}e^{-Q(q)}\Delta(q)^2
                         \dd q_1\cdots\dd q_{n-1}
 =\frac{n!\,\pi^{(n-1)/2}P_n}{\sqrt n},\qquad
 P_n=\prod_{r=1}^{n-1}\frac{r!}{2^r}.
\end{equation}
For completeness, the ordinary $n$-variable Gaussian Vandermonde
integral is $n!\pi^{n/2}P_n$, by the determinant formula and the
Gaussian norms of monic Hermite polynomials. Put $x_i=q_i+m$ with
$\sum q_i=0$. The coordinate Jacobian is $n$, while integration in
$m$ gives $\sqrt{\pi/n}$. Dividing by their product proves
\eqref{eq:mehta}. In particular, replacing the actual independent
coordinate measure by hyperplane surface measure would introduce an
incorrect factor $\sqrt n$.

Write $s_\phi=\sin\phi$ and $\alpha=(n^2-1)/2$. The singular part
of the complete first form factor, on its unique collision chart, has
leading expression
\begin{equation}\label{eq:leadingintegral}
 \frac{2\ii^{1-n^2}B(0)^{n^2/2}}
 {n!(2\pi)^{n-1}2^{n(n-1)}s_\phi^{n^2-1}}
 \int\frac{\Delta(q)^2}{nu+(B(0)/2)Q(q)}
                                   \dd q_1\cdots\dd q_{n-1}.
\end{equation}
This follows directly from
\[
 \gamma(q;u)=\ii\phi-\frac{\ii}{s_\phi}
                  \left(u+\frac{B(0)}2q^2\right)
                  +O(u^2+u q^2+q^4)
\]
and the sine divided differences in the numerator. Every other part of
the fixed-$n$ torus integral is analytic. There are not additional
equal-amplitude collision charts to be omitted or canceled.

Differentiate \eqref{eq:leadingintegral} to the specified order and
scale $q=\sqrt\eps\,v$ on $u=\ii\eps$. The coefficient of $Q$ in
$\gamma$ has positive real part, so the Gaussian/Laplace period is
well-defined with its prescribed branch. The radial tails after the
top derivative have radial integrands $O(R^{-3})$ for odd $n$ and
$O(R^{-2})$ for even $n$; their integrated tails are respectively
$O(R^{-2})$ and $O(R^{-1})$. This gives dominated convergence.
Terms with a higher numerator
degree or an extra $u$ or $Q$ factor in the denominator expansion are
smaller after multiplication by the indicated singularity scale;
splitting at a fixed small chart radius and then at a fixed scaled
radius proves this without an unbounded formal expansion.
Using \eqref{eq:mehta} and the Laplace integral yields
\begin{align}
 \lim_{u\to0,\ u=\ii\eps}u\,\partial_u^kC_n(u)
 &=\frac{2^{1-k}\sqrt{B(0)}P_n n^{k-3/2}}
        {\pi^{(n-1)/2}\sin^{2k}\phi}\ne0,
 &k&=\frac{n^2-1}{2},\quad n\text{ odd},\label{eq:oddfirst}\\
 \lim_{u\to0,\ u=\ii\eps}u^{1/2}\partial_u^kC_n(u)
 &=-\ii\frac{2^{1-\alpha}\sqrt{B(0)}P_n n^{k-1}}
        {\pi^{(n-2)/2}\sin^{n^2-1}\phi}\ne0,
 &k&=\frac{n^2}{2}-1,\quad n\text{ even}.\label{eq:evenfirst}
\end{align}
Here $u^{1/2}=e^{\ii\pi/4}\sqrt\eps$, and $\sqrt{B(0)}$ is fixed
by the beta-period continuation from
$\Re(-\ii B(0)/(2s_\phi))>0$. For even $n$ this sign is determined
by the period; it is not an arbitrary extra square-root branch in the
original integer power of $B$.

For the odd physical HT coefficients the square-root convention is
particularly explicit. On $\OH$, both $\cos y$ and $\cos(\lambda y)$
are positive, and $\sin y$ and $\sin(\lambda y)$ have the same sign.
Consequently
\[
 \Re(T\overline S)=\sinh x\sinh(\lambda x)\cos y\cos(\lambda y)
 +\cosh x\cosh(\lambda x)\sin y\sin(\lambda y)>0
\]
away from the origin, where $B$ has its positive removable value
$\lambda$. Thus $\Re B>0$ throughout the HT lens; its physical square
root is the principal branch. The same is true after axis exchange,
since $\Re(1/B)>0$. This fixes the odd amplitudes on both transfer-axis
charts without an additional sign assumption.

The lower derivatives of the even first term are bounded. For the odd
first term, the $(k-1)$st derivative can be logarithmic, and all lower
derivatives are still $o(u^{-1})$. Together with
\eqref{eq:alltail} this proves that the singularity belongs to the
complete sum, not merely to one term. Multiplication by the analytic
physical prefactor and by $\beta=z/(2J_x)$ preserves it: both are
nonzero at a regular selected point, and the product-rule contributions
from lower derivatives are smaller. A nonramified local change from
$u$ to $z$ also preserves the nonzero top coefficient.

\subsection{All low-temperature signs; the high-temperature distinction}
At a negative endpoint put $q=\pi+q'$. Then
$W=b+u+(-B)(1-\cos q')$. For even $N$, the constraint
$\sum q_i=0\pmod{2\pi}$ remains exactly the same constraint for
$q'_i$, and $B^{N^2/2}=(-B)^{N^2/2}$ because $N^2/2$ is even.
Thus the whole even argument, including its tail and amplitude, applies
to both endpoint signs. The other pair of edges is handled by exchanging
the axes in the symmetric double contour. Rational phase values of the
form above are dense on every regular arc. Section~\ref{sec:geometry}
then supplies the density needed on all of $\partial\OL$.

For odd $N$, the negative endpoint has a different constraint:
$\sum_i(q_i-\pi)$ is an odd multiple of $\pi$. In particular, with
representatives $\delta_i=q_i-\pi\in[-\pi,\pi]$,
\begin{equation}\label{eq:momentummismatch}
 \sum_i\delta_i^2\ge\frac{\pi^2}{N}.
\end{equation}
There is no exact odd all-endpoint collision there. The same Gaussian
numerator bound and \eqref{eq:momentummismatch} imply that the sum of
every fixed parameter derivative is bounded as the negative edge is
approached. This proves existence of all boundary jets, but gives no
positive radius for their Taylor series. The following two sections
resolve this separate infinite-index issue by proving zero radius at
a dense rational-phase family.

% END embedded section: sections/03_regular_edges.tex


% BEGIN embedded section: sections/04_negative_morse.tex
\section{An exact pole normal form at a negative HT edge}
\label{sec:morse}

The odd-particle momentum constraint excludes the all-endpoint
configuration at a negative spectral edge.  The argument in this section
therefore concerns high Taylor coefficients at that edge.  Its estimates
are uniform in the particle number.  In particular, the large order of the
Vandermonde zero is retained in an exact beta integral.

\subsection{Coordinates, the pole equation, and the first quadratic form}

Fix a regular negative endpoint and, if necessary, exchange the two lattice
axes.  In the notation of the angular representation, set
\begin{equation}
 t=\frac{C+B-b}{B},\qquad C=A/S,\qquad
 W(\pi+\delta;t)=b+B(t)\bigl(t-1+\cos\delta\bigr),
 \qquad b=\cos\phi.
 \label{eq:morse-coordinate}
\end{equation}
Here $t=0$ is the boundary point, the pullback is nonramified, and $B(t)$
is holomorphic and nonzero on a fixed disk.  We first take the conjugation
choice
\begin{equation}
 0<\phi<\pi,\qquad \Im B(0)<0,\qquad
 \frac{\phi}{2\pi}=\frac rp,\quad (r,p)=1,\quad p\text{ odd}.
 \label{eq:morse-rational-phase}
\end{equation}
The opposite sign is obtained by conjugation.  Write $s=\sin\phi$ and choose
the analytic branch
\begin{equation}
 g(t,\delta)=\operatorname{arcosh}W(\pi+\delta;t),
 \qquad g(0,0)=\ii\phi.
 \label{eq:morse-g}
\end{equation}
Because $b\in(-1,1)$, this branch is holomorphic on a fixed bidisk in
$(t,\delta)$.  It is even in $\delta$, and $g_t$ is bounded away from zero
after shrinking the bidisk.

Let $n$ tend to infinity through odd multiples of $p$.  In the winding
$+1$ chart of the momentum constraint, introduce
\begin{equation}
 d_n=\frac\pi n,\qquad \tau_n=1-\cos d_n,\qquad
 B_n=B(\tau_n),\qquad
 \delta_i=d_n+v_i,\qquad
 V_n=\left\{v\in\C^n:\sum_i v_i=0\right\}.
 \label{eq:morse-centre}
\end{equation}
The integration measure on the real part of $V_n$ is
$\dd v_1\cdots\dd v_{n-1}$, with $v_n=-\sum_{i<n}v_i$.
Use two distinct quadratic quantities:
\begin{equation}
 Q(v)=\sum_i v_i^2,\qquad q(v)=\sum_i|v_i|^2.
 \label{eq:morse-two-quadratics}
\end{equation}
The first is holomorphic; the second is used only in bounds.
Put
\begin{equation}
 c_n=\frac{\ii B_n}{2s},\qquad
 h_n=\cos d_n-\frac{B_n b}{s^2}\sin^2d_n=1+O(n^{-2}).
 \label{eq:morse-h}
\end{equation}
Direct differentiation at the centre gives
\begin{equation}
 g(\tau_n,d_n)=\ii\phi,\qquad
 g_t(\tau_n,d_n)=-2c_n,\qquad
 \frac12g_{\delta\delta}(\tau_n,d_n)=c_nh_n.
 \label{eq:morse-centre-derivatives}
\end{equation}
Since $n\phi$ is a multiple of $2\pi$, the relevant pole equation is
\begin{equation}
 \Gamma_n(t,v):=\sum_i\bigl(g(t,d_n+v_i)-\ii\phi\bigr)=0.
 \label{eq:morse-pole-equation}
\end{equation}

\begin{lemma}[Uniform implicit pole]
\label{lem:morse-implicit}
There are $r_0>0$, $n_0$, and constants independent of $n$ such that, for
$n\ge n_0$, equation~\eqref{eq:morse-pole-equation} has a unique
permutation-invariant holomorphic solution $t=T_n(v)$ on
$v\in V_n$, $\norm{v}_2<r_0$, with
\begin{align}
 T_n(0)&=\tau_n,&
 |T_n(v)-\tau_n|&\le Cq(v)/n,\label{eq:morse-T-bound}\\
 \norm{\nabla_{V_n}T_n(v)}_2&\le C\sqrt{q(v)}/n,&
 \norm{D^2T_n(v)}&\le C/n.\label{eq:morse-T-derivatives}
\end{align}
All gradients and transposes in holomorphic identities use the complex
bilinear extension of the Euclidean structure on the real zero-sum plane.
\end{lemma}

\begin{proof}
Subtract the constant and linear terms of $g(t,d_n+v_i)$ in $v_i$.
The linear terms sum to zero, the remaining sum is bounded by $Cq(v)$,
and $\Gamma_{n,t}$ has modulus at least $cn$ on a sufficiently small
common neighborhood.  The quantitative implicit-function theorem gives
the first estimate and a single holomorphic solution on a ball whose
radius is independent of $n$.  Uniqueness gives permutation invariance.
Differentiation along $V_n$ gives
\[
 \nabla_{V_n}T_n
 =-\frac{P\bigl(g_\delta(T_n,d_n+v_i)\bigr)_{i=1}^n}
 {\sum_i g_t(T_n,d_n+v_i)},
 \qquad P=I-\frac1n\mathbf1\mathbf1^T.
\]
The numerator has norm at most $C\norm v_2$ and the denominator has
modulus at least $cn$.  A second differentiation gives the stated
Hessian bound.  All analytic derivative bounds come from the original
fixed bidisk; no dimension-dependent Cauchy radius is used.
\end{proof}

\subsection{An exact separable energy identity}

Define the following removable divided differences:
\begin{align}
 a_n(t)&=\frac{g(t,d_n)-\ii\phi}{t-\tau_n},\label{eq:morse-a}\\
 E_n(t,v)&=g(t,d_n+v)-g(t,d_n)-g_\delta(t,d_n)v,\label{eq:morse-E}\\
 H_n(t,v)&=-\frac{2E_n(t,v)}{h_na_n(t)v^2}.
 \label{eq:morse-H}
\end{align}
The function $a_n$ is uniformly bounded away from zero.  Evenness of $g$
in its second argument implies
\begin{align}
 H_n(\tau_n,0)&=1,\label{eq:morse-H-centre}\\
 |H_n(t,v)-1|&\le C\bigl(|t-\tau_n|+|d_n||v|+|v|^2\bigr),
 \label{eq:morse-H-bound}\\
 |\partial_vH_n(t,v)|&\le C(|d_n|+|v|),&
 |\partial_v^2H_n|+|\partial_tH_n|&\le C.
 \label{eq:morse-H-jets}
\end{align}
At the pole, summing~\eqref{eq:morse-E} and using $\sum_i v_i=0$ gives
the exact identity
\begin{equation}
 \boxed{\displaystyle
 T_n(v)-\tau_n=\frac{h_n}{2n}
 \sum_i v_i^2H_n(T_n(v),v_i).}
 \label{eq:morse-energy}
\end{equation}
We next turn its right-hand side into a quadratic form.  Every change of
variable below is holomorphic and permutation equivariant.  In particular,
the construction never divides by $Q(v)$, which has a nontrivial complex
null cone.

\subsection{The projected map and the Vandermonde divisors}

Take the square root of $H_n$ near one and set
\begin{equation}
 U_i(v)=v_i\sqrt{H_n(T_n(v),v_i)},\qquad
 w=G_n(v):=PU(v).
 \label{eq:morse-first-map}
\end{equation}
On a sufficiently small fixed ball in $V_n$,
\begin{equation}
 \norm{DG_n-I}
 \le C\bigl(|d_n|\norm v_\infty+\norm v_\infty^2+q(v)/n\bigr)
 <\frac14.
 \label{eq:morse-first-inverse}
\end{equation}
Thus $G_n$ is uniformly biholomorphic on a smaller fixed ball, and its
inverse has bounded derivative there.  This follows directly by the
uniform inverse-function estimate, or by contraction applied to
$v=w-(G_n(v)-v)$.  Both maps commute with coordinate permutations.

Let $\Delta(v)=\prod_{i<j}(v_i-v_j)$.  The common occurrence of $T_n(v)$
in all components is useful: for every pair,
\begin{equation}
 \frac{U_i-U_j}{v_i-v_j}=1+R_{ij}(v),
 \label{eq:morse-first-divided}
\end{equation}
where the quotient extends holomorphically across $v_i=v_j$ and
\begin{equation}
 |R_{ij}|
 \le C\bigl(|T_n-\tau_n|+|d_n|(|v_i|+|v_j|)
                    +|v_i|^2+|v_j|^2\bigr).
 \label{eq:morse-first-pair-bound}
\end{equation}
Indeed, hold $T_n(v)$ fixed and take the one-variable divided difference
of $z\sqrt{H_n(T_n(v),z)}$.  Projection subtracts a common mean and does
not change component differences.  Sum logarithms of the near-one
factors before applying a norm estimate.  Equations
\eqref{eq:morse-T-bound} and~\eqref{eq:morse-first-pair-bound} give
\begin{align}
 \left|\log\frac{\Delta(w)}{\Delta(v)}\right|
 &\le C\left(n^2|T_n-\tau_n|
         +n|d_n|\sum_i|v_i|+nq(v)\right)\notag\\
 &\le C\bigl(\sqrt{nq(v)}+nq(v)\bigr).
 \label{eq:morse-first-vandermonde}
\end{align}
The quotient here means its holomorphic extension through the coincidence
divisors.

There is also a sharper bound for the Jacobian.  The derivative of
$G_n$ is a compressed diagonal matrix plus one rank-one matrix.  The
diagonal perturbation has trace norm at most
\[
 C\left(|d_n|\sum_i|v_i|+q(v)+n|T_n-\tau_n|\right).
\]
The rank-one trace norm is at most $Cq(v)/n$, by
\eqref{eq:morse-T-derivatives}.  Since its operator norm is below $1/4$,
the convergent matrix logarithm yields
\begin{equation}
 |\log\det_{V_n}DG_n|
 \le C\bigl(|d_n|\sqrt{nq(v)}+q(v)\bigr).
 \label{eq:morse-first-jacobian}
\end{equation}
The determinant on $V_n$ is also the change-of-variable determinant in
the constrained coordinates $v_1,\ldots,v_{n-1}$.

\subsection{The mean correction without division by a quadratic form}

Express $v$ as a function of $w$ and put
\begin{equation}
 \mu_n(w)=\frac1n\sum_iU_i(v(w)).
 \label{eq:morse-mean}
\end{equation}
The common linear coefficient of $U_i$ drops out of this sum.  More
explicitly, at fixed $t$ the expansion has the form
\[
 U_i=a_0(t)v_i+a_1(t)v_i^2+a_2(t,v_i)v_i^3,
 \qquad a_1(t)=O(d_n),
\]
with uniformly bounded analytic coefficients.  Consequently, with
$q=q(w)$,
\begin{align}
 |\mu_n(w)|&\le\frac Cn(|d_n|+\sqrt q)q,\label{eq:morse-mean-bound}\\
 \norm{\nabla\mu_n(w)}_2&\le\frac Cn(|d_n|+\sqrt q)\sqrt q,\label{eq:morse-mean-gradient}\\
 \norm{D^2\mu_n(w)}&\le\frac Cn(|d_n|+\sqrt q).
 \label{eq:morse-mean-hessian}
\end{align}
To obtain these bounds, differentiate the displayed expansion and use
$\sum v_i=0$, the estimates for $T_n$, and the bounded inverse of $G_n$.
Terms from differentiating $T_n$ have an additional factor $1/n$.
The second derivative of $G_n^{-1}$ is bounded by
$C(|d_n|+\norm w_2)$, so composition preserves
\eqref{eq:morse-mean-hessian}.

The holomorphic Hadamard formula on $V_n$ gives
\begin{equation}
 \mu_n(w)=w^T\ell_n(w),\qquad
 \ell_n(w)=\int_0^1\nabla\mu_n(\theta w)\dd\theta\in V_n.
 \label{eq:morse-hadamard}
\end{equation}
The field $\ell_n$ is permutation equivariant and satisfies
\begin{equation}
 \norm{\ell_n}_2\le\frac Cn(|d_n|+\sqrt q)\sqrt q,
 \qquad
 \norm{D\ell_n}\le\frac Cn(|d_n|+\sqrt q).
 \label{eq:morse-ell-bounds}
\end{equation}
In particular,
\begin{equation}
 \ell_{n,i}-\ell_{n,j}=(w_i-w_j)L_{ij}(w),
 \qquad |L_{ij}|\le\frac Cn(|d_n|+\sqrt q).
 \label{eq:morse-ell-divided}
\end{equation}
For the quantitative assertion, compare $w$ to the vector obtained by
exchanging its $i,j$ entries.  Equivariance identifies the corresponding
output difference, and the derivative bound along the straight segment
gives the displayed inequality.  Holomorphic divisibility follows by the
one-variable divided-difference formula, including on the hyperplane of
coincident entries.

Now define
\begin{align}
 R_n(w)&=\left(I+n\ell_n(w)\ell_n(w)^T\right)^{1/2}w
        =w+a_n^\sharp(w)\mu_n(w)\ell_n(w),\label{eq:morse-rank-map}\\
 a_n^\sharp(w)&=\frac n{1+\sqrt{1+n\ell_n(w)^T\ell_n(w)}}.
 \label{eq:morse-rank-coefficient}
\end{align}
These are analytic matrix square roots near the identity.  Transposes in
these identities are bilinear transposes, not Hermitian adjoints.  The
matrix perturbation has operator norm at most
$C(|d_n|+\sqrt q)^2q/n$, so $R_n$ is uniformly biholomorphic on a
smaller fixed ball.  Exact multiplication gives
\begin{equation}
 Q(R_n(w))=Q(w)+n\mu_n(w)^2=\sum_iU_i(v)^2.
 \label{eq:morse-rank-energy}
\end{equation}
This identity is valid on the complex null cone as well.

By~\eqref{eq:morse-ell-divided}, the relative correction to each
difference $w_i-w_j$ under $R_n$ is at most
$C(|d_n|+\sqrt q)^2q/n$.  Its summed logarithmic cost over all pairs
is therefore at most $Cnq$ on the fixed ball.  Differentiating
\eqref{eq:morse-rank-map} gives a trace-norm bound
$C(|d_n|+\sqrt q)^2q$ for $DR_n-I$, and hence the same bound for its
logarithmic Jacobian.  In particular, this step introduces no factor
of the form $n^2\norm v_\infty^2$ into the final density estimate.

\begin{lemma}[Exact symmetric normal form]
\label{lem:negative-morse}
On a common fixed ball in $V_n$, for all sufficiently large admissible
$n$, the permutation-equivariant biholomorphism
\begin{equation}
 \Psi_n(v)=\sqrt{h_n}\,R_n(G_n(v)),
 \qquad D\Psi_n(0)=\sqrt{h_n}I,
 \label{eq:morse-final-map}
\end{equation}
with the square root near one, satisfies
\begin{equation}
 \boxed{T_n(v)=\tau_n+\frac{Q(\Psi_n(v))}{2n}.}
 \label{eq:morse-exact-normal-form}
\end{equation}
It preserves every coincidence divisor.  After extracting the constant
factor $h_n^{-\alpha_n}$, where $\alpha_n=(n^2-1)/2$, the
Jacobian--Vandermonde ratio has a holomorphic logarithm bounded by
\begin{equation}
 C\bigl(\sqrt{nq}+nq\bigr).
 \label{eq:morse-geometric-density}
\end{equation}
The constants are independent of $n$.
\end{lemma}

\begin{proof}
Combine~\eqref{eq:morse-energy} with~\eqref{eq:morse-rank-energy}.
The divided-difference and Jacobian bounds above give
\eqref{eq:morse-geometric-density}, also for the inverse map after
reducing the common ball.  The degree of $\Delta^2$ is $n(n-1)$ and
the dimension is $n-1$.  Scaling by $\sqrt{h_n}$ therefore extracts
$h_n^{-[n(n-1)+(n-1)]/2}=h_n^{-\alpha_n}$.
\end{proof}

\subsection{The complete form-factor density at the pole}

In this chart include the constrained measure normalization in the
numerator:
\begin{equation}
 F_n(t,v)=\frac{B(t)^{n^2/2}}{n!(2\pi)^{n-1}}
 \prod_i\frac1{\sinh g(t,d_n+v_i)}
 \prod_{i<j}\left[
 \frac{\sin((v_i-v_j)/2)}
 {\sinh((g(t,d_n+v_i)+g(t,d_n+v_j))/2)}\right]^2.
 \label{eq:morse-complete-numerator}
\end{equation}
The square-root sign in the odd power of $B$ is the physical one.
The global kernel is $\coth(\Gamma_n/2)$ on this resonant chart.
Its pole contributes the Taylor coefficient density
\begin{equation}
 \frac{\rho_n(v)}{T_n(v)^{j+1}},\qquad
 \rho_n(v)=-\frac{2F_n(T_n(v),v)}
 {\Gamma_{n,t}(T_n(v),v)}.
 \label{eq:morse-residue-density}
\end{equation}
The minus sign comes from expanding $(t-T_n)^{-1}$ at $t=0$.
The quotient of $\rho_n(v)$ by $\Delta(v)^2$ is holomorphic and
nonzero.  At the centre its value is
\begin{equation}
 \rho_{n,0}=-\frac{2B_n^{\alpha_n-1/2}}
 {n\,n!(2\pi)^{n-1}2^{n(n-1)}s^{2\alpha_n}}.
 \label{eq:morse-centre-density}
\end{equation}

\begin{lemma}[Uniform full-density estimate]
\label{lem:morse-density}
In the coordinates $w=\Psi_n(v)$, the form
$\rho_n(v)\dd v_1\cdots\dd v_{n-1}$ is
\begin{equation}
 \rho_{n,0}h_n^{-\alpha_n}\Delta(w)^2
 \bigl(1+\mathcal E_n(w)\bigr)
 \dd w_1\cdots\dd w_{n-1},
 \label{eq:morse-transformed-density}
\end{equation}
where, uniformly in $n$ on a common fixed ball,
\begin{equation}
 |\mathcal E_n(w)|\le
 \exp\!\left\{C\bigl(\sqrt{nq(w)}+nq(w)\bigr)\right\}-1.
 \label{eq:morse-density-error}
\end{equation}
\end{lemma}

\begin{proof}
Divide the entire density by its centre value and by $\Delta^2$ before
estimating it.  Every remaining analytic factor is nonzero on a common
small ball, so its logarithm is available there.  The power $n^2/2$
of $B(T_n)$ costs $Cn^2|T_n-\tau_n|$.  For the one-body phases,
\[
 \sum_i|g(T_n,d_n+v_i)-\ii\phi|
 \le C\left(|d_n|\sum_i|v_i|+q(v)+n|T_n-\tau_n|\right).
\]
The sum over all pair-denominator logarithms costs at most $Cn$ times
this last quantity.  The sine divided differences cost
$C\sum_{i<j}|v_i-v_j|^2\le Cnq(v)$.  One-body denominators have the
smaller single-index bound.  Finally, cancellation of the linear term
under $\sum v_i=0$ gives
\[
 \frac{\Gamma_{n,t}(T_n(v),v)}{n g_t(\tau_n,d_n)}
 =1+O(q(v)/n).
\]
Thus the logarithmic cost of the complete residue density is at most
\[
 C\left(n^2|T_n-\tau_n|+n|d_n|\sum_i|v_i|+nq(v)\right)
 \le C\bigl(\sqrt{nq(v)}+nq(v)\bigr).
\]
Combine this estimate with Lemma~\ref{lem:negative-morse}.  No
uncontrolled constant raised to $n^2$ is used in this comparison with
the centre value.
\end{proof}

\subsection{The real cycle, its deformation, and the boundary remainder}

\begin{lemma}[Cycle and analytic remainder]
\label{lem:morse-cycle-remainder}
Localize to the winding $+1$ real chart
$\sum_i\delta_i^2<\sigma$, with fixed sufficiently small $\sigma>0$.
Its domain in $V_n\cap\R^n$ is the actual ball
$Q(v)<\sigma-\pi^2/n$.  One may deform its pole integral to a real
ball in the exact coordinates of Lemma~\ref{lem:negative-morse}.
The boundary contribution and the pole-subtracted analytic remainder
have Taylor coefficients bounded by
\begin{equation}
 (C n)^j\exp\{-\kappa n^2+C n\log(n+1)\},
 \qquad j\ge0,\quad n\ge n_0,
 \label{eq:morse-boundary-remainder}
\end{equation}
where $C,\kappa>0$ are independent of $n,j$.  The same holds for
the winding $-1$ chart.
\end{lemma}

\begin{proof}
First choose a fixed ball on which $\Psi_n$ and its inverse estimates
hold.  Uniform near-identity bounds show that its image contains a
fixed smaller ball.  Next choose the real integration ball sufficiently
small that its image and the straight homotopy between it and the
original ball lie in this common inverse-map domain.  This order of
choices ensures that the transformed form is defined on the actual
deformation chain.  A fixed smaller ball can be used instead of the
full ball $Q<\sigma-\pi^2/n$; the intervening shell obeys the same
estimate below.

For small enough radius and large enough $n$,
\[
 \norm{\Psi_n(v)-v}_2\le\eta\norm v_2,\qquad \eta<1/4.
\]
If $v$ is real and $0\le u\le1$, then
\begin{equation}
 \Re Q\bigl((1-u)v+u\Psi_n(v)\bigr)
 \ge(1-2\eta-\eta^2)Q(v)>0.
 \label{eq:morse-homotopy-positive}
\end{equation}
Consequently the denominator $\tau_n+Q(w)/(2n)$ cannot vanish on
this homotopy.  Stokes' theorem for the holomorphic top-degree form
changes the integral only by its boundary tube.  On that tube
$\Re Q(w)\ge c\sigma$, and hence
\[
 \left|\tau_n+\frac{Q(w)}{2n}\right|\ge\frac{c_\sigma}{n}.
\]

The particle bound on this genuinely complex chain follows from the
exact matched-Pfaffian identity.  All angles remain in a fixed small
complex neighborhood of $\pi$, and all spectral roots remain close
to $\exp(-\ii\phi)$.  Every $z$-Schur factor therefore has a common
bound $\rho_0<1$, while every canceled complete pair is bounded.
Expanding the odd augmented Pfaffian gives
$\exp\{-\kappa n^2+O(n\log(n+1))\}$ after the one-body and factorial
factors are included.  No assertion that the roots have modulus at
most one on the complex deformation is required.  The factor
$\Gamma_{n,t}^{-1}$ costs at most $C/n$.  The maps, inverse maps,
sphere volumes, and boundary parameterization cost at most
$\exp(O(n\log(n+1)))$.  This proves
\eqref{eq:morse-boundary-remainder} for the boundary tube and shell.

For completeness, the pole subtraction has the same scale.  Let $C_T$
bound $n|T_n(v)|/\sigma$ on the local ball. After shrinking $\sigma$,
choose a fixed $c_{\rm pole}>2C_T\sigma$ with
$C_\Gamma(c_{\rm pole}+\sigma)<\pi$, where $C_\Gamma$ is a uniform
bound for the phase increment. Thus $T_n(v)$ lies inside
$|t|<c_{\rm pole}/(2n)$. On $|t|<c_{\rm pole}/n$,
$\Gamma_n$ has only the displayed zero,
$\Gamma_n/(t-T_n)$ has modulus at least $c_1n$, and
$|\Gamma_n|<\pi$.  Thus
$\coth(\Gamma_n/2)-2/\Gamma_n$ is analytic there.  Divide
$2F_n$ by $\Gamma_n/(t-T_n)$, subtract its value at $T_n$, and
divide the remaining analytic difference by $t-T_n$.
The resulting analytic remainder has the same matched-Pfaffian
bound, with polynomial factors in $n$ absorbed into
$\exp(O(n\log(n+1)))$.  Cauchy's estimate on a disk of radius
$c_{\rm pole}/(4n)$ gives~\eqref{eq:morse-boundary-remainder}. This
local subtraction disk is chosen independently of the smaller exterior
disk in Section~\ref{sec:jets}. Reflection
$\delta\mapsto-\delta$ treats the other winding.
\end{proof}

\subsection{The exact beta period and the uniform two-winding formula}

Write
\begin{equation}
 \alpha_n=\frac{n^2-1}{2},\qquad
 P_n=\prod_{r=1}^{n-1}\frac{r!}{2^r}.
 \label{eq:morse-alpha-P}
\end{equation}
The constrained Gaussian period is
\begin{equation}
 \mathcal G_n:=\int_{V_n\cap\R^n}e^{-Q(v)}\Delta(v)^2
                   \dd v_1\cdots\dd v_{n-1}
 =\frac{n!\pi^{(n-1)/2}P_n}{\sqrt n}.
 \label{eq:morse-gaussian}
\end{equation}
One way to verify it is to expand the Vandermonde determinant in monic
Hermite polynomials, whose squared norms are
$\sqrt\pi\,2^{-r}r!$.  The resulting unconstrained integral is
$n!\pi^{n/2}P_n$.  Separating the mean coordinate has Jacobian $n$
in the displayed constrained coordinates and contributes
$\sqrt{\pi n}$, giving~\eqref{eq:morse-gaussian}.
Homogeneity and the beta integral now give, when $j+1>\alpha_n$,
\begin{align}
 &\int_{V_n\cap\R^n}
 \frac{\Delta(w)^2\dd w_1\cdots\dd w_{n-1}}
 {(\tau_n+Q(w)/(2n))^{j+1}}\notag\\
 &\hspace{8mm}=
 \mathcal G_n(2n)^{\alpha_n}\tau_n^{\alpha_n-j-1}
 \frac{\Gamma(j+1-\alpha_n)}{\Gamma(j+1)}.
 \label{eq:morse-beta-period}
\end{align}
The gamma ratio is retained exactly.  Replacing it by
$j^{-\alpha_n}$ at fixed $n$ would have an error with logarithm of
order $\alpha_n^2/j$, which is not small in the regime used below.

\begin{lemma}[Uniform beta-density error]
\label{lem:negative-beta-uniform}
The relative effect of the density in
Lemma~\ref{lem:morse-density}, after the cycle deformation, is
\begin{equation}
 O\left(\frac n{\sqrt j}+\frac{n^2}{j}\right)
 \label{eq:morse-beta-relative-error}
\end{equation}
uniformly as $n\to\infty$ and $n^2/j\to0$.  Extension of the
unperturbed radial integral from the local ball to the whole plane
has an exponentially small relative error in this regime.
\end{lemma}

\begin{proof}
Under the normalized radial beta measure set
$x=Q/(2n\tau_n)$.  Its density is proportional to
$x^{\alpha_n-1}(1+x)^{-j-1}$, and
\[
 nQ\asymp x,\qquad
 \mathbb E x=\frac{\alpha_n}{j-\alpha_n}.
\]
The comparison constants are uniform because $2n^2\tau_n$ stays
between positive constants.  On $0\le x\le\delta$ with fixed
$\delta>0$, the right side of~\eqref{eq:morse-density-error} has
expectation at most
$C\sqrt{\alpha_n/j}+C\alpha_n/j$, using
$\mathbb E\sqrt x\le\sqrt{\mathbb E x}$.

Here is an explicit tail estimate that also controls the exponential
density bound.  The truncated ball has $x\le Ln$.  For
$\delta\le x\le Ln$ and $j\gg n^2$,
\[
 C\sqrt x+Cx\le\tfrac14j\log(1+x).
\]
If
\[
 Z(\alpha,J)=\int_0^\infty x^{\alpha-1}(1+x)^{-J-1}\dd x,
\]
the normalized remaining integral is at most
\begin{equation}
 (1+\delta)^{-j/4}\frac{Z(\alpha_n,j/2)}{Z(\alpha_n,j)}
 \le (1+\delta)^{-j/4}4^{\alpha_n},
 \label{eq:morse-beta-exp-tail}
\end{equation}
provided $\alpha_n\le j/4$.  In the present problem $\alpha_n$ is
an integer, and the last ratio is exactly
\[
 \prod_{r=0}^{\alpha_n-1}\frac{j-r}{j/2-r}\le4^{\alpha_n}.
\]
Thus~\eqref{eq:morse-beta-exp-tail} decreases exponentially in $j$
uniformly when $\alpha_n/j\to0$.  The same estimate without the
exponential density factor bounds the omitted tail of the unperturbed
beta integral.  This proves~\eqref{eq:morse-beta-relative-error}.
\end{proof}

\begin{proposition}[Uniform local Taylor coefficients at a negative edge]
\label{prop:negative-local-jets}
Let $a^{\mathrm{loc}}_{n,j}$ denote the Taylor coefficient of $C_n$
from the two local windings $+1$ and $-1$ described above.  Uniformly
for odd multiples $n$ of the fixed odd $p$, as $n\to\infty$ and
$n^2/j\to0$,
\begin{equation}
 a^{\mathrm{loc}}_{n,j}
 =M(n,j)\left[1+O\left(\frac n{\sqrt j}+\frac{n^2}{j}\right)\right]
 +R_{n,j},
 \label{eq:morse-final-relative}
\end{equation}
where
\begin{equation}
 \boxed{\displaystyle
 M(n,j)=
 -\frac{2^{2-\alpha_n}B_n^{\alpha_n-1/2}P_n n^{\alpha_n-3/2}}
 {\pi^{(n-1)/2}s^{2\alpha_n}}
 h_n^{-\alpha_n}\tau_n^{\alpha_n-j-1}
 \frac{\Gamma(j+1-\alpha_n)}{\Gamma(j+1)},}
 \label{eq:morse-model}
\end{equation}
and
\begin{equation}
 |R_{n,j}|\le (Cn)^j
 \exp\{-\kappa n^2+C n\log(n+1)\}.
 \label{eq:morse-final-remainder}
\end{equation}
All constants are independent of $n,j$ in this regime.
\end{proposition}

\begin{proof}
Multiply the centre density~\eqref{eq:morse-centre-density}, the
factor $h_n^{-\alpha_n}$, and the exact period
\eqref{eq:morse-beta-period}.  Reflection $\delta\mapsto-\delta$
maps the two winding charts to one another; their coefficients are
equal and add, giving the factor two in~\eqref{eq:morse-model}.
Lemma~\ref{lem:negative-beta-uniform} gives the relative error, and
Lemma~\ref{lem:morse-cycle-remainder} gives
\eqref{eq:morse-final-remainder}.
\end{proof}

\begin{remark}
The phase of $B_n^{\alpha_n-1/2}$ varies quadratically with $n$.
Formula~\eqref{eq:morse-model} does not assert positivity or a common
phase for different particle indices.  The comparison of indices in
the next section is needed for noncancellation of the complete odd
sum.  The uniform regime here includes, in particular,
$\log j\le n\le\sqrt j/(\log j)^{1/4}$; the constants in
\eqref{eq:morse-final-relative} do not change across that interval.
\end{remark}

% END embedded section: sections/04_negative_morse.tex


% BEGIN embedded section: sections/05_negative_jets.tex
\section{Negative high-temperature edges: boundary jets and the complete odd sum}
\label{sec:jets}

The absence of an odd-particle endpoint resonance at a negative spectral
edge does not imply continuation of the infinite odd sum. We prove its
nonanalyticity directly from its boundary Taylor coefficients. This avoids
assuming a common continuation contour through the nearby interior
finite-particle singularities.

\subsection{Setting and the existence of the boundary jets}

Fix a regular negative endpoint. After complex conjugation if necessary,
the local analytic data can be written
\begin{equation}\label{eq:jets_setting}
 W(t,\pi+\delta)=b+B(t)\bigl[t-(1-\cos\delta)\bigr],\qquad
 b=\cos\phi,\quad \phi=\frac{2\pi r}{p}\in(0,\pi),
\end{equation}
where $p\ge3$ is odd, $(r,p)=1$, $B$ is holomorphic and nonzero,
$\Im B(0)<0$, and $t$ is a nonramified local temperature coordinate.
The physical inward direction is $t=-\eps$, $\eps>0$. The branch
$\cosh\gamma(q;t)=W(t,q)$ has
$\gamma(\pi;0)=\ii\phi$. All assertions below concern this branch.
At a regular negative edge, the earlier root construction gives a common
branchfree parameter neighborhood for the entire real angular circle.

For clarity, write the complete coefficient as
\begin{equation}\label{eq:jets_full_coefficient}
 C_n(t)=\int_{\mathcal T_n}H_n(t,q)
        \coth\!\left(\frac{\Gamma_n(t,q)}2\right)\dd q,
 \qquad \Gamma_n=\sum_{i=1}^n\gamma(q_i;t),
\end{equation}
where $\mathcal T_n$ is the zero-total-momentum torus,
$q_n=-\sum_{i<n}q_i$ modulo $2\pi$, the measure $\dd q$ is the
ordinary measure in the $n-1$ independent angles, and
\begin{equation}\label{eq:jets_numerator}
 H_n(t,q)=\frac{B(t)^{n^2/2}}{n!(2\pi)^{n-1}}
 \prod_{i=1}^n\frac1{\sinh\gamma(q_i;t)}
 \prod_{i<k}\left[
  \frac{\sin((q_i-q_k)/2)}
       {\sinh((\gamma(q_i;t)+\gamma(q_k;t))/2)}
 \right]^2.
\end{equation}
The branch of $B^{1/2}$ is the one in the physical representation. The
normalization in \eqref{eq:jets_numerator} is retained throughout this
section. The physical susceptibility differs from the complete odd sum
by the analytic prefactor already identified in the physical
representation; that prefactor is nonzero at a regular endpoint.

Set
\[
 S(q)=\sum_{i=1}^n\dist(q_i,\pi)^2.
\]
The spectral geometry and the matched Schur--Pfaffian estimate imply,
with constants independent of $n$,
\begin{align}
 \Re\gamma(q;0)&\ge c_0\dist(q,\pi)^2,
 &|\gamma(q;0)-\ii\phi|&\le C\dist(q,\pi)^2,\label{eq:jets_attenuation0}\\
 \sup_{|t|<r_0,\ q\in\mathcal T_n}|H_n(t,q)|
 &\le E_n,
 &E_n&=\exp\{-\kappa n^2+C n\log(n+1)+C\}.
 \label{eq:jets_gaussian_num}
\end{align}
Shrinking $r_0$ is allowed. In particular, the strict Schur contraction
used in \eqref{eq:jets_gaussian_num} persists on this fixed complex
parameter disk; it does not require $|e^{-\gamma}|\le1$ at every
complex parameter in the disk.

For odd $n$, choose the shortest lifts $q_i=\pi+\delta_i$ modulo
$2\pi$. Momentum conservation gives
$\sum_i\delta_i=w\pi$ with $w$ odd. Therefore
\begin{equation}\label{eq:jets_momentum_minimum}
 S(q)\ge\frac{\pi^2}{n}.
\end{equation}
This is the source of an $n^{-2}$ coefficientwise analytic disk, rather
than a common disk for the entire particle sum.

\begin{lemma}[Coefficientwise disks and boundary derivatives]
\label{lem:jets_boundary_derivatives}
There is $c>0$ such that each odd coefficient is holomorphic on
$|t|<c/n^2$, agrees there with \eqref{eq:jets_full_coefficient}, and
obeys
\begin{equation}\label{eq:jets_small_disk}
 \sup_{|t|<c/n^2}|C_n(t)|\le CnE_n,
 \qquad |[t^j]C_n(t)|\le (Cn^2)^j CnE_n.
\end{equation}
The complete odd sum on $t<0$ has derivatives of every fixed order at
$t=0$ from that side, and its normalized boundary jets are
\begin{equation}\label{eq:jets_boundary_coefficients}
 a_j=\lim_{\eps\downarrow0}\frac1{j!}
       \partial_t^j\sum_{n\ \mathrm{odd}}C_n(-\eps)
     =\sum_{n\ \mathrm{odd}}[t^j]C_n(t).
\end{equation}
The sum on the right converges absolutely for each fixed $j$.
\end{lemma}

\begin{proof}
The parameter derivative of each one-particle phase is bounded on a
fixed disk. Thus, using \eqref{eq:jets_attenuation0} and
\eqref{eq:jets_momentum_minimum}, a sufficiently small perturbation
$|t|<c/n^2$ preserves $\Re\Gamma_n\ge c'/n$. The modulus of the
hyperbolic cotangent is then at most $Cn$. This proves the first bound;
Cauchy's estimate proves the second.

For the inward real approach, \eqref{eq:jets_setting} gives
\[
 \Re\gamma(q;-\eps)\ge
 c\bigl(\eps+\dist(q,\pi)^2\bigr).
\]
Indeed $\Im W=(-\Im B(-\eps))(\eps+1-\cos\delta)$ is positive,
and writing $\gamma=v+\ii\psi$ expresses it as
$\sinh v\sin\psi$, with $\psi$ in a compact subinterval of
$(0,\pi)$. A disk centered at $-\eps$ of radius
$c(\eps+n^{-2})$ consequently preserves the same denominator gap,
up to a fixed factor. Cauchy's estimate there bounds the $j$th
derivative by $j!(Cn^2)^j CnE_n$, uniformly for small $\eps$.
This is summable for every fixed $j$. Differentiation, the boundary
limit, and the particle sum therefore commute by dominated
convergence. This argument proves existence of all boundary jets; it
does not assert a common holomorphic disk for their sum.
\end{proof}

\subsection{Removing all analytic remainders}

Choose $\sigma>0$ sufficiently small, and then a fixed
$c_{\rm out}>0$ such that
$C_\Gamma c_{\rm out}<c_0\sigma/2$ and
$C_\Gamma(\sigma+c_{\rm out})<\pi/p$.
Here $C_\Gamma$ bounds the phase increments below. On
$|t|<c_{\rm out}/n$,
\begin{equation}\label{eq:jets_c_over_n_estimate}
 \Re\Gamma_n\ge c_0S-C_\Gamma c_{\rm out},\qquad
 |\Gamma_n-n\ii\phi|\le C_\Gamma(S+c_{\rm out}).
\end{equation}
If $S\ge\sigma$, the first estimate supplies a fixed positive real
part by this choice of $c_{\rm out}$. If $S<\sigma$ and $p\nmid n$, the
second estimate supplies a fixed phase gap: the distance from
$n\phi$ to $2\pi\Z$ is at least $2\pi/p$. Thus the whole
coefficient for $p\nmid n$, and the part with $S\ge\sigma$ for any
$n$, are holomorphic on this $c_{\rm out}/n$ disk with norm at most $CE_n$.
The angular cutoff is fixed and independent of $t$.

For $p\mid n$, put $D_n=\Gamma_n-n\ii\phi$. On the inner region
$S<\sigma$, $|D_n|<\pi$ on the same disk if the constants are
sufficiently small. The identity
\begin{equation}\label{eq:jets_coth_subtraction}
 \coth(D_n/2)=\frac2{D_n}+R(D_n),
 \qquad R \text{ holomorphic on } |D_n|<2\pi,
\end{equation}
places the $R$ part in the same analytic remainder class.

We also subtract the exact simple parameter pole, not only its
quadratic approximation. In the inner region use the unique small
lifts $q_i=\pi+\delta_i$ and put $v_i=1-\cos\delta_i$. Let
\[
 d_n(t,q)=\frac1nD_n(t,q),\qquad
 a_0=\frac{B(0)}{\sinh(\ii\phi)}\ne0.
\]
On a sufficiently small fixed parameter disk,
\[
 \partial_t d_n(t,q)=a_0+O(|t|+\sigma),
\]
uniformly in $n$ and in $S<\sigma$. Reducing the disk and $\sigma$
makes this derivative lie within $|a_0|/4$ of $a_0$. The quantitative
inverse-function estimate gives a unique zero
$t=T_n(q)=O(S/n)$ and the factorization
\begin{equation}\label{eq:jets_exact_pole_factor}
 D_n(t,q)=n(t-T_n(q))A_n(t,q),\qquad
 A_n(t,q)=\int_0^1\partial_t d_n
       (T_n(q)+u(t-T_n(q)),q)\dd u.
\end{equation}
The factor $A_n$ is analytic and uniformly bounded away from zero on
a smaller fixed disk containing every such $T_n(q)$. The proof can
also be obtained directly by Rouch\'e's theorem on a circle of radius
$CS/n$: $|d_n(0,q)|\le CS/n$, while the derivative differs uniformly
from the nonzero constant $a_0$.

Set $F_n=H_n/A_n$. The exact identity
\begin{equation}\label{eq:jets_divided_difference}
 \frac{2H_n(t,q)}{D_n(t,q)}=
 \frac{2F_n(T_n(q),q)}{n(t-T_n(q))}
 +\frac2n\frac{F_n(t,q)-F_n(T_n(q),q)}{t-T_n(q)}
\end{equation}
has a holomorphic second term, including at $t=T_n(q)$. A
fixed-disk Cauchy bound on $\partial_tF_n$ bounds this divided
difference by $CE_n$. In particular no lower bound on the distance
between the pole and a chosen integration contour is being assumed.

\begin{lemma}[Exact reduction to the inner pole integrals]
\label{lem:jets_remainder}
For every $j\ge0$ the boundary coefficient in
\eqref{eq:jets_boundary_coefficients} is the sum of the exact inner
pole coefficients for odd $n$ divisible by $p$, together with a
remainder $r_j$ satisfying
\begin{equation}\label{eq:jets_remainder_bound}
 |r_j|\le \exp\left\{\frac12j\log(j+1)+Cj+C\right\}.
\end{equation}
For a resonant $n$, the inner pole coefficient is exactly
\begin{equation}\label{eq:jets_exact_residue_coefficient}
 a_{j,n}^{\mathrm{pole}}=
 -\int_{S<\sigma}
  \frac{2H_n(T_n(q),q)}{\partial_t\Gamma_n(T_n(q),q)}
  T_n(q)^{-j-1}\dd q.
\end{equation}
\end{lemma}

\begin{proof}
All removed terms are bounded on a disk $c_{\rm rem}/n$, where
$c_{\rm rem}$ is a fixed sufficiently small fraction of
$\min(c_{\rm pole},c_{\rm out})$. The exact pole subtraction is
constructed on its own larger disk before restriction; the pole need
not lie in the exterior disk. The preceding decompositions are identities on the $c/n^2$ disk
where the original integral is holomorphic. Taking its Taylor
coefficients gives \eqref{eq:jets_exact_residue_coefficient}; the
minus sign is the coefficient of $(t-T)^{-1}$ at $t=0$. Every
removed term is analytic on $|t|<c_{\rm rem}/n$ with norm bounded by $CE_n$.
Their total coefficient bound is
\[
 \sum_{n\ge1}(Cn)^j E_n.
\]
Absorb $Cn\log(n+1)$ into $(\kappa/2)n^2+C'$, and optimize
$j\log n-(\kappa/2)n^2$. This gives
\eqref{eq:jets_remainder_bound}, including summability of the tail.
\end{proof}

\subsection{Uniform bounds for the winding sectors}

The exact pole is quantitatively close to the average
$\bar v=n^{-1}\sum_i v_i$. At $t=\bar v$, the linear terms in the
analytic expansion of
$\operatorname{arcosh}(b+B(\bar v)(\bar v-v_i))$ sum to zero.
The same uniform inverse-function estimate as above consequently gives
\begin{equation}\label{eq:jets_pole_lower}
 T_n(q)=\bar v+O\bigl(\operatorname{Var}(v)\bigr),\qquad
 |T_n(q)|\ge (1-C\sigma)\frac{S(q)}{2n}.
\end{equation}
Here $\operatorname{Var}(v)=n^{-1}\sum_i(v_i-\bar v)^2\le
C\sigma\bar v$, and
$\bar v\ge (1-C\sigma)S/(2n)$. The second estimate concerns the
modulus of a generally complex pole; it does not assert that the
pole is real.

Write $s=\sin\phi>0$, $\alpha_n=(n^2-1)/2$, and
\begin{equation}\label{eq:jets_H0}
 H_{0,n}=\frac{|B(0)|^{n^2/2}}
 {n!(2\pi)^{n-1}2^{n(n-1)}s^{n^2}}.
\end{equation}
The complete numerator, evaluated at the pole, satisfies
\begin{equation}\label{eq:jets_numerator_vandermonde}
 |H_n(T_n(q),q)|\le H_{0,n}\Delta(\delta)^2 e^{CnS},
 \qquad \Delta(\delta)=\prod_{i<k}(\delta_i-\delta_k).
\end{equation}
To check the constant and the dependence on $n$, divide every sine
numerator by $(\delta_i-\delta_k)/2$ and every hyperbolic
denominator by $\sinh(\ii\phi)$. Analytic logarithms of the
resulting near-one ratios are bounded by
$C(\delta_i-\delta_k)^2$ and
$C(|T_n|+\delta_i^2+\delta_k^2)$, respectively. Their sums are
$O(nS+n^2|T_n|)=O(nS)$. The one-body factors cost $O(S)$.
Finally the logarithm of
$B(T_n)^{n^2/2}/B(0)^{n^2/2}$ costs $O(n^2|T_n|)$.
This proves \eqref{eq:jets_numerator_vandermonde} with the exact
scale \eqref{eq:jets_H0}, rather than an unspecified constant raised
to the power $n^2$.

In winding sector $w$, where $w$ is odd, put
\[
 \delta_i=\frac{w\pi}{n}+\eta_i,\quad
 \sum_i\eta_i=0,\quad Q=\sum_i\eta_i^2,\quad
 m_{n,w}=\frac{\pi^2w^2}{n}.
\]
Then $S=Q+m_{n,w}$ and $\Delta(\delta)=\Delta(\eta)$.
There are at most $C\sqrt n$ sectors with nonempty $S<\sigma$.
The constrained Gaussian integral, with the same independent-angle
measure as above, is
\begin{equation}\label{eq:jets_Mehta}
 I_n=\int_{\R^{n-1}}e^{-Q}\Delta(\eta)^2\dd\eta
 =\frac{\pi^{(n-1)/2}2^{-n(n-1)/2}}{\sqrt n}
   \prod_{r=1}^n r!.
\end{equation}
For example, separating the center of mass in the ordinary Gaussian
Vandermonde integral gives a Jacobian $n$ and a center Gaussian
$\sqrt{\pi/n}$, whose product explains the constrained
$1/\sqrt n$ factor. Homogeneity therefore gives, whenever
$j+1>\alpha_n$,
\begin{equation}\label{eq:jets_beta_integral}
 \int_{\R^{n-1}}\frac{\Delta(\eta)^2}{(Q+m)^{j+1}}\dd\eta
 =I_n m^{\alpha_n-j-1}
       \frac{\Gamma(j+1-\alpha_n)}{\Gamma(j+1)}.
\end{equation}

\begin{lemma}[Higher windings and large particle numbers]
\label{lem:jets_windings_large_n}
For $j+1>\alpha_n$, the exact pole coefficient in winding $w$
satisfies
\begin{align}
 |a_{j,n,w}^{\mathrm{pole}}|
 &\le \frac{C H_{0,n}}n e^{Cn\sigma}
 \left(\frac{2n}{1-C\sigma}\right)^{j+1}
 I_n m_{n,w}^{\alpha_n-j-1}
 \frac{\Gamma(j+1-\alpha_n)}{\Gamma(j+1)}.
 \label{eq:jets_winding_bound}
\end{align}
For every resonant odd $n$ and every $j\ge0$, without a restriction
involving $\alpha_n$,
\begin{equation}\label{eq:jets_large_n_bound}
 \log\sum_w|a_{j,n,w}^{\mathrm{pole}}|
 \le 2j\log n-\frac12n^2\log n+Cn^2+Cj+C.
\end{equation}
Empty sectors are understood to contribute zero.
\end{lemma}

\begin{proof}
The residue denominator has modulus at least $cn$. Combine this
with \eqref{eq:jets_pole_lower} and
\eqref{eq:jets_numerator_vandermonde}, replace $S$ in its exponential
by $\sigma$, and extend the positive radial majorant to the whole
real zero-sum plane. Equation \eqref{eq:jets_beta_integral} proves
\eqref{eq:jets_winding_bound}. Its exact dependence on $w$ relative
to the same expression with $w=1$ is
$|w|^{-2(j+1-\alpha_n)}$.

For the second estimate, use the simpler bound $|T_n|\ge c/n^2$
and the exact ball mass
\[
 \int_{Q<\sigma}\Delta(\eta)^2\dd\eta
 =I_n\frac{\sigma^{\alpha_n}}{\Gamma(\alpha_n+1)}.
\]
The logarithm of this ball mass is
$-\tfrac12n^2\log n+O(n^2)$; this follows by applying Stirling's
inequalities to \eqref{eq:jets_Mehta} and to
$\Gamma(\alpha_n+1)$. Multiplication by $H_{0,n}$ costs $O(n^2)$.
The number of windings and the remaining polynomial factors are
absorbed into the displayed error. This proves
\eqref{eq:jets_large_n_bound} for arbitrarily large $n$.
\end{proof}

\subsection{The exact first-winding model}

Let
\begin{equation}\label{eq:jets_model_parameters}
 \tau_n=1-\cos(\pi/n),\quad B_n=B(\tau_n),\quad
 h_n=\cos(\pi/n)-\frac{B_n b}{s^2}\sin^2(\pi/n),\quad
 P_n=\prod_{r=1}^{n-1}r!\,2^{-r}.
\end{equation}
In particular $h_n=1+O(n^{-2})$. The exact symmetric pole normal
form and cycle estimate of Proposition~\ref{prop:negative-local-jets}
in Section~\ref{sec:morse} give, uniformly
as $n\to\infty$ through odd multiples of $p$ and $n^2/j\to0$,
\begin{equation}\label{eq:jets_uniform_model}
 a_{j,n,1}^{\mathrm{pole}}+a_{j,n,-1}^{\mathrm{pole}}
 =M(n,j)(1+e_{n,j})+b_{n,j},
\end{equation}
where
\begin{equation}\label{eq:jets_model}
 M(n,j)=-\frac{2^{2-\alpha_n}B_n^{\alpha_n-1/2}P_n
                     n^{\alpha_n-3/2}}
                   {\pi^{(n-1)/2}s^{2\alpha_n}}
 h_n^{-\alpha_n}\tau_n^{\alpha_n-j-1}
 \frac{\Gamma(j+1-\alpha_n)}{\Gamma(j+1)},
\end{equation}
and
\begin{equation}\label{eq:jets_model_errors}
 |e_{n,j}|\le C\left(\frac n{\sqrt j}+\frac{n^2}j\right),\qquad
 |b_{n,j}|\le(Cn)^j
       \exp\{-\kappa' n^2+C'n\log(n+1)+C'\}.
\end{equation}
The same estimate applies to the actual inner ball of radius squared
$\sigma-\pi^2/n$, since these radii stay between fixed positive
constants for large $n$. Alternatively its outer shell belongs to a
$c/n$ analytic remainder by the preceding argument.

All constants in \eqref{eq:jets_model} can be checked directly.
At the collision $\gamma_i=\ii\phi$, the numerator divided by
$\Delta^2$ is
\[
 -\frac{\ii B_n^{n^2/2}}
       {n!(2\pi)^{n-1}2^{n(n-1)}s^{n^2}},
 \qquad \Gamma_t=-\frac{\ii nB_n}s.
\]
Multiplication by $-2/\Gamma_t$, followed by the beta integral in
the exact normal form, gives half of \eqref{eq:jets_model}.
Reflection $\delta\mapsto-\delta$ preserves the full integrand and
exchanges the first two winding sectors, giving the displayed
factor of two. The factor $B_n^{\alpha_n-1/2}$ retains its quadratic
particle-dependent phase. No positivity across particle indices is
being assumed.

\subsection{A discrete concavity estimate with controlled errors}

Write $L(n,j)=\log|M(n,j)|$ and set $d=2p$, the step between
successive allowed odd multiples of $p$. Define
\begin{equation}\label{eq:jets_c0}
 c_*=-\frac34-\frac32\log2+\log\pi
       +\frac12\log|B(0)|-\log s.
\end{equation}
Elementary Stirling estimates in the exact formula give
\begin{equation}\label{eq:jets_L_expansion}
 L(n,j)=2j\log n-\frac{n^2}2\log j+c_*n^2
       -j\log(\pi^2/2)
       +O(n\log n+n^4/j+j/n^2),
\end{equation}
when $n\to\infty$ and $n^2/j\to0$.
We do not differentiate this big-O term. The following discrete
estimate is obtained from the exact factors instead.

\begin{lemma}[Exact-factor second differences]
\label{lem:jets_discrete_concavity}
Uniformly when $n\to\infty$ and $n^2/j\to0$,
\begin{equation}\label{eq:jets_second_difference}
 L(n+d,j)-2L(n,j)+L(n-d,j)
 =d^2\left(-\frac{2j}{n^2}-\log j+2c_*\right)
  +O_d(n^2/j+j/n^4+n^{-2}).
\end{equation}
In particular this difference is at most $-c d^2\log j$ throughout
the interval
\begin{equation}\label{eq:jets_middle_window}
 \log j<n<\frac{\sqrt j}{(\log j)^{1/4}},
\end{equation}
for all sufficiently large $j$, whenever the adjacent indices are
also in that interval.
\end{lemma}

\begin{proof}
The exact quotient of the factorial products gives
\begin{align*}
 \Delta_d^2\log P_n
 &=\sum_{k=0}^{d-1}\sum_{\ell=1}^d
       \log(n-d+k+\ell)-d^2\log2\\
 &=d^2(\log n-\log2)+O_d(n^{-2}).
\end{align*}
The offsets in the double sum add to zero, so their first-order
$1/n$ terms cancel.

For integer $j>\alpha_{n+d}$, the gamma quotient can be written
\[
 \log\frac{\Gamma(j+1-\alpha_n)}{\Gamma(j+1)}
 =-\alpha_n\log j+
   \sum_{m=0}^{\alpha_n-1}-\log(1-m/j).
\]
The two adjacent increments of $\alpha_n$ have lengths
$dn+d^2/2$ and $dn-d^2/2$. Their $d^2$ excess summands cost
$O_d(n^2/j)$. Pairing the remaining summands costs the same amount,
because the function $-\log(1-x/j)$ has derivative at most
$1/(j-\alpha_{n+d})$ on the required interval. Thus this factor has
second difference $-d^2\log j+O_d(n^2/j)$.

Finally
\[
 \log\tau_n=\log(\pi^2/2)-2\log n
             -\frac{\pi^2}{12n^2}+O(n^{-4}),
\]
while $B(\tau_n)=B(0)(1+O(n^{-2}))$ and $h_n=1+O(n^{-2})$
have convergent expansions in $1/n$ near zero. Computing the
second differences of these explicit terms in
\eqref{eq:jets_model} gives
\eqref{eq:jets_second_difference}, including the error $j/n^4$
from the exact value of $\tau_n$.

In \eqref{eq:jets_middle_window}, the error $j/n^4$ is smaller than
$j/n^2$ by a factor at most $(\log j)^{-2}$, and $n^2/j$ tends
uniformly to zero. The claimed negative bound follows. This proves
concavity across the whole intermediate range, not just a fixed
relative neighborhood of its maximum.
\end{proof}

\subsection{Selecting an integer order with one dominating index}

Fix a large allowed index $n$ and let
$F_n(j)=L(n+d,j)-L(n-d,j)$, initially with real
$j\asymp n^2\log n$. The two gamma arguments differ by the
integer $2dn$, so differentiation of their exact quotient gives
\begin{align}
 F_n'(j)
 &=\log\frac{\tau_{n-d}}{\tau_{n+d}}
    -\sum_{m=1}^{2dn}\frac1{j+1-\alpha_{n-d}-m}\notag\\
 &=\frac{4d}{n}-\frac{2dn}{j}
    +O_d(n^{-3}+n^3/j^2)>0.
 \label{eq:jets_crossing_derivative}
\end{align}
The exact first difference, with the ordinary Stirling bounds for
the finitely many factorials in each quotient, also gives
\begin{equation}\label{eq:jets_crossing_expansion}
 \frac{F_n(j)}{2d}=\frac{2j}{n}-n\log j+2c_*n
       +O_d(\log n+n^3/j+j/n^3).
\end{equation}
At
\[
 j=n^2\bigl(\log n+\tfrac12\log\log n-c_*\pm\eps\bigr)
\]
this expression has opposite signs for every fixed $\eps>0$ and
all sufficiently large $n$. Consequently there is a unique
crossing in this range. Round it to the nearest integer and denote
the result by $j_n$. Then
\begin{equation}\label{eq:jets_selected_j}
 j_n=n^2\bigl(\log n+\tfrac12\log\log n-c_*+o(1)\bigr),
 \qquad F_n(j_n)=O_d(1/n).
\end{equation}
The selected $n$ lies strictly inside
\eqref{eq:jets_middle_window} for these $j_n$.

Let $\ell_k=L(n+kd,j_n)$ for the allowed indices in that window.
Lemma~\ref{lem:jets_discrete_concavity} says that successive
differences of $\ell_k$ decrease by at least $c d^2\log j_n$.
Equation \eqref{eq:jets_selected_j} makes the two adjacent values
equal up to $O_d(1/n)$. It follows first that
$\ell_0-\ell_{\pm1}\ge c'd^2\log j_n$, and then, by summing
the decreasing first differences, that
\[
 \ell_k-\ell_0\le -c''d^2 k^2\log j_n+O_d(|k|/n)
 \qquad(k\ne0).
\]
In particular
\begin{equation}\label{eq:jets_single_index}
 \sum_{\substack{m\ne n,\ p\mid m,\ m\ \mathrm{odd}\\
                  m\text{ in }\eqref{eq:jets_middle_window}}}
 |M(m,j_n)|
 \le |M(n,j_n)|\,O(j_n^{-c'''d^2}).
\end{equation}
The varying complex phases are immaterial to this modulus
inequality. An $o(1)$ relative error in the local model is
sufficient: the error of the largest index is relative to its own
modulus, while the errors at other indices multiply already
smaller moduli.

\subsection{Completing the all-index comparison}

For $j\to\infty$, split the odd particle indices at
\[
 n_{\mathrm{lo}}=\log j,\qquad
 n_{\mathrm{hi}}=\frac{\sqrt j}{(\log j)^{1/4}}.
\]
The low range, including every nonasymptotic index, is bounded by
Lemma~\ref{lem:jets_boundary_derivatives}:
\begin{equation}\label{eq:jets_low_range}
 \sum_{n\le n_{\mathrm{lo}}}|[t^j]C_n|
 \le\exp\{2j\log\log j+O(j)\}.
\end{equation}
All $c/n$ analytic remainders and the normal-form boundary errors
have the smaller scale
$\exp\{\tfrac12j\log j+O(j)\}$ by
\eqref{eq:jets_remainder_bound} and
\eqref{eq:jets_model_errors}.

For $n\ge n_{\mathrm{hi}}$, use the ball-mass bound
\eqref{eq:jets_large_n_bound}. In the variable $x=n^2$ its
controlling exponent is
\[
 j\log x-\tfrac14x\log x+Cx+Cj,
\]
which is decreasing for every $x\ge j/\sqrt{\log j}$ when $j$
is sufficiently large. At the threshold it equals at most
\begin{equation}\label{eq:jets_high_range}
 j\log j-\tfrac12j\log\log j-\tfrac14j\sqrt{\log j}
       +o(j\sqrt{\log j}).
\end{equation}
Its subsequent decrease makes the full infinite tail summable,
with the same leading bound. This estimate handles all windings in
the high range and requires no beta integral outside its
convergence range.

In the intermediate range, the error in
\eqref{eq:jets_uniform_model} is uniformly
$O((\log j)^{-1/4})=o(1)$. It remains to compare the higher odd
windings to that model. Let $M_0(n,j)$ be the modulus of
\eqref{eq:jets_model} with $B_n$ replaced by $B(0)$, $h_n$ by
$1$, and $\tau_n$ by $\pi^2/(2n^2)$. The exact formulas give
\begin{equation}\label{eq:jets_baseline_ratio}
 \left|\log\frac{M_0(n,j)}{|M(n,j)|}\right|
 \le C(j/n^2+1)=o(j)
\end{equation}
uniformly in this entire range. Equation
\eqref{eq:jets_winding_bound} is bounded by $CM_0(n,j)$ times
\[
 e^{Cn\sigma}(1-C\sigma)^{-j-1}
 |w|^{-2(j+1-\alpha_n)}.
\]
Here $n=o(j)$ and $\alpha_n/j=o(1)$ uniformly. Choose the fixed
$\sigma$ so small that $-\log(1-C\sigma)<\tfrac12\log3$.
Summing over all odd $|w|\ge3$ then gives
\begin{equation}\label{eq:jets_higher_windings_negligible}
 \sum_{|w|\ge3}|a_{j,n,w}^{\mathrm{pole}}|
 \le |M(n,j)|e^{-c j}
\end{equation}
for some $c>0$ and all sufficiently large $j$, uniformly over the
intermediate range. Thus no unexamined winding or particle range
remains between the preceding estimates.

At the selected sequence \eqref{eq:jets_selected_j}, the exact
model and \eqref{eq:jets_L_expansion} yield
\begin{align}
 \log|M(n,j_n)|
 &=j_n\log j_n-j_n\log\log j_n
   +j_n\bigl(\log(4/\pi^2)-1\bigr)+o(j_n),
 \label{eq:jets_main_growth}\\
 |M(n,j_n)|^{1/j_n}
 &\sim\frac{4j_n}{e\pi^2\log j_n}.
 \notag
\end{align}
For this calculation, the exact gamma ratio differs from
$-\alpha_n\log j_n$ in its logarithm by $O(\alpha_n^2/j_n)$,
which is $o(j_n)$ at the selected scale. It was retained exactly
in the finite differences and the relative local estimate, where
discarding it would not have been legitimate.

The bounds \eqref{eq:jets_low_range} and
\eqref{eq:jets_high_range}, and all analytic remainders, are
$o(|M(n,j_n)|)$ by \eqref{eq:jets_main_growth}. The intermediate
higher windings have the same property by
\eqref{eq:jets_higher_windings_negligible} and
\eqref{eq:jets_single_index}. Finally the uniform relative error
in \eqref{eq:jets_uniform_model}, summed over that range, is at
most $o(1)\sum_m|M(m,j_n)|=o(|M(n,j_n)|)$. Therefore the
complete boundary coefficient satisfies
\begin{equation}\label{eq:jets_complete_noncancellation}
 a_{j_n}=M(n,j_n)(1+o(1)).
\end{equation}
This proves noncancellation in the infinite odd-particle sum even
though the model phases vary with $n$.

\begin{theorem}[A regular rational negative edge is singular for the full odd sum]
\label{thm:jets_negative_edge}
Under the physical-branch and nonramified-coordinate hypotheses
of \eqref{eq:jets_setting}, the complete odd sum has all inward
boundary derivatives at $t=0$, but the Taylor series of those
boundary derivatives has radius of convergence zero. More
precisely, along the integer sequence \eqref{eq:jets_selected_j},
\[
 |a_{j_n}|^{1/j_n}\sim\frac{4j_n}{e\pi^2\log j_n}\longrightarrow\infty.
\]
Consequently the complete high-temperature physical germ has no
holomorphic continuation across this regular negative edge
point. Multiplication by its nonzero analytic physical
normalization preserves this conclusion.
\end{theorem}

\begin{proof}
The jets exist and equal the absolutely convergent coefficient
sums by Lemma~\ref{lem:jets_boundary_derivatives}. Equations
\eqref{eq:jets_main_growth} and
\eqref{eq:jets_complete_noncancellation} give their stated
growth and zero Taylor radius. If there were a holomorphic
continuation across the point, its Taylor coefficients would be
these same boundary jets and would satisfy a positive-radius
Cauchy bound, a contradiction. The physical normalization can
be divided out in a neighborhood because it is analytic and
nonzero there. The use of the nonramified coordinate $t$ makes
this a statement about the fixed-ratio one-variable physical
germ itself.
\end{proof}


% END embedded section: sections/05_negative_jets.tex


% BEGIN embedded section: sections/06_closure.tex
\section{Closure, exceptional points, and the isotropic interface}\label{sec:closure}

\subsection{Completing the unequal-coupling theorem}
We now trace the two physical germs to their full boundaries. The
representation in Section~\ref{sec:physical} gives holomorphic functions
on the explicit, independently determined domains $\OH$ and $\OL$.
Their equality with the two physical susceptibilities holds on the
appropriate positive real intervals. The identity theorem fixes the
continued physical germs on the entire corresponding connected domains.

For each fixed $\lambda\ne1$, the exceptional set in
Section~\ref{sec:geometry} is locally discrete. Every other boundary
point has a real analytic edge arc with exactly one endpoint cosine,
a strictly interior second cosine, a nonzero real-linear determinant,
and a nonramified holomorphic edge coordinate. Restricting that
coordinate to the arc gives an open real interval. Cosines of primitive
odd-prime rational angles are dense in every such interval.

On every regular LT arc, Section~\ref{sec:regular} proves a genuine
singularity of the entire even-particle sum at each point of this dense
family, including both endpoint signs and both transfer-axis choices.
The first singular derivative has a nonzero coefficient and every
higher even index is included in the absolute bounded tail. On every
regular positive HT arc, the odd counterpart proves the analogous
whole-sum singularity. On every regular negative HT arc,
Sections~\ref{sec:morse}--\ref{sec:jets} show that all fixed boundary
jets exist, but a sequence of their Taylor coefficients satisfies
\begin{equation}\label{eq:boundarygrowthclosure}
 |a_{j_n}|^{1/j_n}\sim\frac{4j_n}{e\pi^2\log j_n}\longrightarrow\infty.
\end{equation}
These coefficients are those of the complete odd sum. The Gaussian
bound permits taking the boundary limit through the sum for each fixed
derivative. If a holomorphic extension existed, its derivatives would
have these limits and its Taylor series would have positive radius,
contradicting \eqref{eq:boundarygrowthclosure}. Division by the local
analytic nonzero physical prefactor proves the same assertion for
$\XH$, and multiplication by the nonzero $\beta$ gives $\chi_H$.

Finally suppose a holomorphic extension existed through any boundary
point, including an exceptional point. A small disk of extension contains
regular arcs of the same physical boundary, and hence contains a nearby
point of the appropriate dense true-singularity family. Restricting the
extension to that point contradicts its already proved nonremovability.
The complete graph $\partial\OL$ and the complete closed curve
$\partial\OH$ therefore are natural boundaries. This proves
the unequal-coupling part of Theorem~\ref{thm:aniso}.

This argument never identifies all of $\cS_\lambda$ with the singular
set of susceptibility. For unequal couplings $\cS_\lambda$ has interior;
the barriers are the explicitly determined boundaries of its two
physical complementary components. Nor is any existence or density
claim about singular differential equations used as an infinite-sum
noncancellation theorem.

\subsection{All ratios and degeneracies in the new argument}
The only ratio excluded from the strict anisotropic-edge mechanism is
$\lambda=1$, where the real-linear determinant degenerates along the
whole isotropic spectral curve. No rationality, algebraicity,
Diophantine bound, or proximity to one is imposed on the other ratios.
For $\lambda\ne1$, the exceptional boundary points consist of the
locally discrete collinear points, corners, elementary zeros or poles
in a chosen chart, and critical points of the edge coordinate. They
are included in the final barrier by density, not declared removable.
Simultaneous imaginary-axis cosine zeros occur exactly for positive
odd-over-odd rational ratios; their local geometry is analyzed in
Section~\ref{sec:geometry} and is likewise included. The near-real-critical
opening of the two spectral edges is nonzero for every unequal positive
ratio, so there is no omitted interval of weak anisotropy.

The negative HT proof uses $B(t)$ only as a fixed holomorphic function
near a regular point. Its exact resonant parameters
$1-\cos(\pi/n)$ and strict single-index estimates do not require
arithmetic control of the temperature coordinate or of $\lambda$.
This is why highly approximable irrational ratios are covered without
an implicit generic-ratio qualification.

\subsection{Exact reuse on the isotropic slice}
The following interface is an identity of the actual formulas, rather
than an inference from the labels attached to previous projects. Set
$\lambda=1$ and $s=\sinh z$. Then $S=T=s$, $p=s^2$, $B=1$, and
\[
 D(x,y;z)=s\left[s+s^{-1}
                  -\tfrac12(x+x^{-1})-\tfrac12(y+y^{-1})\right].
\]
On the positive LT interval $\sqrt p=s$, so every factor
$\sqrt p/D$ in \eqref{eq:doublecorr} is exactly the normalized
isotropic factor in \cite{LTsource}. The even coefficients and the
full-site identity consequently agree term by term, with
$\xi_L=(1-s^{-4})^{1/4}$. In the HT interval the full angular coefficient
in \eqref{eq:fullangular} has $W=s+s^{-1}-\cos q$, and its prefactor
is exactly $(1-s^4)^{1/4}/s$. This is the literal normalized odd sum
in \cite{HTsource,HTreport}, not a diagonal susceptibility or a dual
low-temperature observable.
In the notation of that HT manuscript its auxiliary exterior variable
is $\widehat t=1/s$; hence our $C_N$ at $s$ is exactly its
$C_N(\widehat t)$, and the reported real condition $\widehat t>1$
is exactly the present physical interval $0<s<1$.

The isotropic results to be reused at this interface are precisely:
\begin{enumerate}
\item the exterior connected full-bulk result in
\cite{LTsource}, Corollary labeled \texttt{thm:nb}, derived there from
its whole-form-factor theorem \texttt{thm:first} and all-index theorem
\texttt{thm:tail};
\item the normalized HT analytic germ and unit-circle obstruction of
\cite{HTsource}, Theorem 1.1, with the normalized analytic endpoint
reported in \cite{HTreport}. Its physical identification requires the
real full-bulk form-factor equality, absolute onsite summability and
onsite normalization, all supplied by Section~\ref{sec:physical} at
$B=1$.
\end{enumerate}
The later report's asserted kernel validation is kept at that literal
scope. It did not derive Gibbs states from the Hamiltonian, does not
establish an anisotropic theorem, and was not replayed in this campaign.
The low-temperature manuscript's candidate label is not a mathematical
axiom. Its actual first/tail argument was inspected at the minimum
load-bearing reuse cut. In its notation, at its dense prime family,
$N_0=2p$ and $k=N_0^2/2-1$, it proves
\[
 T_{N_0}^{(k)}((1+\eps)s_*)
 =2L_*\eps^{-1/2}+o(\eps^{-1/2}),\qquad L_*\ne0,
\]
and, for every fixed $0\le j\le k$,
\[
 \sum_{\substack{N>N_0\\N\ {\rm even}}}
 |T_N^{(j)}((1+\eps)s_*)|=o(\eps^{-1/2}).
\]
The inspected arguments include the entire first-coefficient torus
complement, its two noncanceling local charts, the weighted Stokes
current, the unfactored high-order Vandermonde/Lie estimates, and every
particle window. No open premise at this normalized analytic scope
remained in that inspection. Their literal resulting function is
$1-M^2+2M^2\sum_{N\ {\rm even},N\ge2}T_N$, which is exactly
\eqref{eq:globalLT} on the isotropic slice by \eqref{eq:allsites}.
The HT source similarly supplies a nonzero first odd coefficient and
the entire differentiated odd tail, on both signs of the boundary real
part; its normalized function is exactly \eqref{eq:globalHT}.
These are reused analytic theorems about explicitly defined functions.
The real Gibbs-state and uniform-field identifications are supplied
here by Section~\ref{sec:physical}, not inferred from either analytic
theorem's label. The detailed exact-scope reconstruction and source
identities are retained in the dependency and verification records.

The geometry on this slice is exact:
\[
 \cS_1=\{z:|\sinh z|=1\},\qquad
 H_1(y)=\operatorname{arsinh}(\cos y),\quad |y|<\pi/2,
 \qquad L_1(y)=\operatorname{arsinh}(|\cos y|).
\]
The map $z\mapsto\sinh z$ maps the central HT cell biholomorphically
onto $|s|<1$. It maps $\OL$ as the universal cover of $|s|>1$;
its $2\pi\ii$ period is therefore a covering statement, not a claim
that the energy-field response $\chi=\beta\mathcal X$ is globally
single-valued as a function of $s$. At all nonaxis boundary points the
map is locally biholomorphic. The boundary critical points with
$\cosh z=0$ are limits of such points. Dense isotropic obstructions
therefore pull back to every point of the two boundaries, including
those critical points. The factor $\beta=z/(2J_x)$ is nonzero on the
selected boundary family.

\begin{proof}[Completion of the proof of Theorem~\ref{thm:aniso}]
The unequal-ratio part was proved at the start of this section. At the equal ratio,
apply the two reused normalized results through the exact identities
and covering maps above. No limit in $\lambda$ and no phase duality
is used. This proves the theorem for every finite positive ratio on
both physical temperature sides.
\end{proof}

\begin{corollary}[Inverse-temperature barriers]\label{cor:alllambda}
For each fixed $J_x>0$, the barriers in the $\beta$ plane are exactly
$\partial\OH/(2J_x)$ and $\partial\OL/(2J_x)$ for their respective
physical germs. At $\lambda=1$ they are the correctly interpreted
pullbacks of $|\sinh(2\beta J)|=1$, with the low-temperature covering
specified above.
\end{corollary}

\subsection{Verification and computational scope}
The analytic arguments carry all infinite-index and dense-boundary
claims. The accompanying code checks exact finite polynomial identities
for the spectral geometry and Gaussian periods, tests the full nonlinear
three-particle singular coefficient by direct differentiated integration,
and tests the new symmetric normal form, including a complex null-cone
configuration. Separate exact-model computations test the finite-difference
and single-index mechanism at large indices and on a fixed $\lambda=2$
slice. The latter are computations of the displayed beta model; they
are not high-dimensional quadratures of the complete susceptibility.
All such diagnostics retain their actual precision and scope in their
JSON evidence files. Compilation verifies presentation only.

The campaign uses genuinely separate mathematical worker and audit
contexts. These constitute internal adversarial checking, not external
human review, publication-priority clearance, or a new Lean validation.
The final disposition of each dependency and each destructive finding
is given in the accompanying verification record. The three assigned
new-mathematics cuts and their affected repair checks were completed;
the root then reconstructed their joint physical and all-ratio scope.
No publication-priority clearance, external human validation, or new
Lean replay is claimed.

% END embedded section: sections/06_closure.tex

\begin{thebibliography}{99}
\bibitem{BL} A. I. Bugrij and O. Lisovyy, Correlation function of the two-dimensional
Ising model on a finite lattice. II, arXiv:0708.3643 (2007), especially
equations (21)--(25), (30)--(39).
\bibitem{IL} N. Iorgov and O. Lisovyy, Ising correlations and elliptic determinants,
arXiv:1012.2856 (2010), especially Theorem 6.1, Remark 6.2,
and the vacuum specializations.
\bibitem{ONGP} W. P. Orrick, B. G. Nickel, A. J. Guttmann and J. H. H. Perk,
The susceptibility of the square lattice Ising model: New developments,
arXiv:cond-mat/0103074 (2001), especially equations (1.9)--(1.11),
(3.13), (3.19), (3.24), and Section 3.2.
\bibitem{BHM} S. Boukraa, S. Hassani and J.-M. Maillard,
Holonomic functions of several complex variables and singularities of
anisotropic Ising n-fold integrals, arXiv:1207.1784v2 (2012),
especially the unsquared anisotropic Nickel equation.
\bibitem{Raoufi} A. Raoufi, Translation-invariant Gibbs states of Ising model:
General setting, arXiv:1710.07608 (2017), Theorem 1.
\bibitem{CGN} F. Camia, C. Garban and C. M. Newman, The Ising magnetization
exponent on $\mathbb Z^2$ is $1/15$, arXiv:1205.6612v4 (2013),
Theorem 1.5 (the GHS inequality for arbitrary nonnegative interactions
and site fields) and Corollary 1.6 (pinned plus spins).
\bibitem{LTsource} D. Lin, A candidate natural-boundary proof for the bulk
square-lattice Ising susceptibility, supplied frozen version 0.1-rc6,
17 October 2026, file \texttt{LT\_Ising\_v0.1-rc6\_source\_frozen.tex}.
The exact theorem and its physical premises are identified in Section~\ref{sec:closure}.
\bibitem{HTsource} Supplied high-temperature Ising manuscript HT-1.0,
file \texttt{HT\_Ising\_HT-1.0\_source.tex}, Theorem 1.1 and its normalized
odd-particle representation.
\bibitem{HTreport} Supplied 18 October 2026 HT Lean formalization status JSON
and final report. This campaign inspected their stated theorem and premise
scope; it did not replay the reported build.
\end{thebibliography}
\end{document}
